\chapter{Design Patterns}
\label{chap:designpatterns}
\vspace{10.5cm}
\textbf{Design Patterns} describe how to solve recurring design problems 
to design flexible and reusable object-oriented software, that is, 
objects that are easier to implement, change, test, and reuse.

\newpage
%%%
%%% CREATIONAL PATTERNS %%%
%%%
\section{Creational Patterns}

%%% %%% FACTORY METHOD %%% %%% %%%
\subsection{Factory Method}
\begin{table}[h!t]
    \begin{center}
        \begin{tabularx}{\textwidth}{|l|X|}
            \hline
            \textbf{What for} & \highlightit{creating objects} without 
                                having to specify the exact class 
                                of the \\ 
                              & object that will be created\\
            \hline
            \textbf{How to} & a \highlightit{factory method} is used 
                              \highlightit{instead of calling a constructor}:
                              it can be \\
                            & either \highlightit{specified in an interface} 
                              and implemented by child classes, \\
                            & or \highlight{implemented in a base class} 
                              and optionally overridden by \\
                            & derived classes\\

            \hline
            \textbf{When needed} & $\bullet$ exact types and dependencies 
                                   of the objects are unknown \\
                                 & \quad beforehand \\
                                 & $\bullet$ internal components of 
                                   a library or framework need to be 
                                   extended\\
                                 & \quad by users\\ 
                                 & $\bullet$ system resources need to 
                                 be saved by reusing existing objects. \\
            \hline
        \end{tabularx}
    \end{center}
\end{table}

\vspace{-1cm}
\toclessnparagraph{Structure}
\vspace{-1.5cm}
\begin{table}[h!t]
    \begin{center}
    \begingroup
    \begin{figure}[H]
        \begin{center}
        %~ \begin{center}
    \begin{tikzpicture}[show background grid]
            
        \begin{abstractclass}[text width=3cm]{ABCBuilder}{0,0}
            \operationtt{+ reset()}
            \operationtt{+ buildStepA()}
            %~ \operationtt{+ buildStepB()}
            \operationtt{+ buildStepZ()}
        \end{abstractclass}
        \begin{class}[text width=4.5cm]{BuilderA}{-3,-4.5}
            \implement{ABCBuilder}
            \attributett{result: ProductA}
            \operationtt{+ reset()}
            \operationtt{+ buildStepA()}
            %~ \operationtt{+ buildStepB()}
            \operationtt{+ buildStepZ()}
            \operationtt{+ getResult(): ProductA}
        \end{class}
        \begin{class}[text width=4.5cm]{BuilderB}{3,-4.5}
            \implement{ABCBuilder}
            \attributett{result: ProductB}
            \operationtt{+ reset()}
            \operationtt{+ buildStepA()}
            %~ \operationtt{+ buildStepB()}
            \operationtt{+ buildStepZ()}
            \operationtt{+ getResult(): ProductB}
        \end{class}


        \begin{class}[text width=4.5cm]{Director}{5,-0.5}
            \attributett{- builder: Builder}
            \operationtt{+ Director(builder)}
            \operationtt{+ changeBuilder(builder)}
            \operationtt{+ make(type)}
        \end{class}

    \end{tikzpicture}
\end{center}

            %~ \begin{center}
\centering
    \tikzfig{factory_method_design_pattern}
\begin{tikzpicture}%[show background grid]
        \begin{class}[text width=5cm]{Creator}{0,0}
            \attributett{...}
            \operationtt{+ someOperation()}
            \operationtt{+ factoryMethod(): Product}
        \end{class}
        
        %~ \begin{class}[text width=5cm]{CreatorA}{-2.5,-5.5}
        \begin{class}[text width=5cm]{CreatorA}{0,-4}
            \inherit{Creator}
            \attributett{...}
            \operationtt{+ factoryMethod(): ProductA}
        \end{class}

        %~ \begin{class}[text width=5cm]{CreatorB}{2.5,-5.5}
            %~ \inherit{Creator}
            %~ \attributett{...}
            %~ \operationtt{+ createProduct(): Product}
        %~ \end{class}
        
        \begin{interface}[text width=3cm]{Product}{8,-0.2}
            \operationtt{+ doStuff()}
        \end{interface}
        
        \begin{class}[text width=3.3cm]{ProductA}{8,-3.7}
        \implement{Product}\end{class}

        %~ \begin{class}[text width=3.3cm]{ConcreteProductB}{10,-4}
        %~ \implement{Product}\end{class}
        
        %~ \unidirectional
        
        \draw [ - >] ( Creator.east) ++ (0,-0.15) -- 
        %~ node [ above , sloped , black ]{$ < <$ instantiate $ > >$} 
        ( Product.west);
        \draw [ umlcd style dashed line , - >] ( CreatorA.east) ++ (0,-0.5) 
        -- ++(2,0) node [ above , sloped , black ]{$ < <$ create $ > >$} 
        -- ++(2,0) -|
        %~ node [ above , sloped , black ]{$ < <$ instantiate $ > >$} 
        ( ProductA) ;
        %~ \aggregation{Abstraction}{}{}{Implementation}
    \end{tikzpicture}
%~ \end{center}

        \end{center}
    \end{figure}
    \endgroup
    \vspace{-1.5cm}
        \begin{tabularx}{\textwidth}{|l|X|}
            \hline
            \highlightbf{ABCProduct} & interface common to all 
                                       the objects. \\
            \hline
            \highlightbf{ConcreteProduct} & different implementations 
                                             of the product interface. \\
            \hline
            \highlightbf{Creator} & base class which declares 
                                       the factory method: \\
                                     & \small{\, $\bullet$ \highlightbf{factory method} 
                                        - abstract or returning a default product;} \\
                                     & \small{\, $\bullet$ \highlightbf{returned type} 
                                       - must match the interface.} \\
                                     & It can have some \highlightit{core business logic} 
                                       which can be splitted from the 
                                       product creation. \\
            \hline
            \highlightbf{CreatorA},  & concrete classes that \highlightit{override the base factory method} \\
            \highlightbf{CreatorB}   & to return the different types of concrete products. \\
            \hline
        \end{tabularx}
        \caption{Factory Method Class Diagram.}
        \label{tab:factory-method-class-diagram}
    \end{center}
\end{table}

\newpage
\toclessnparagraph{Pros and Cons}
\begin{table}[h!t]
    \begin{tabularx}{\textwidth}{|l|X|}
        \hline
        \textbf{Pros} & $\bullet$\, \highlightit{decoupling} of the 
                        creation of concrete objects from the core 
                        business logic \\
                      & \quad of the Creator; \\
                      & $\bullet$\, \highlightit{Single Respobsibility Principle}\\
                      & $\bullet$\, \highlightit{Open/Close Principle} \\
        \hline
        \textbf{Cons} & code gets more complex \\
        \hline      
    \end{tabularx}
\end{table}
    
%~ \toclessnparagraph{Relation with other Design Patterns}
%~ \begin{table}[h!t]
    %~ \begin{tabularx}{\textwidth}{|l|X|}
        %~ \hline
        %~ \highlightbf{Abstract Factory} & can be based on a set of 
                                         %~ \highlightit{Factory Method}s\\
        %~ \hline
        %~ \highlightbf{Iterator} & can be used along with \highlightit{Factory Method} 
                                 %~ to let collection subclasses have 
                                 %~ their own compatible iterators\\
        %~ \hline
        %~ \highlightbf{Prototype} & isn't based on inheritance but 
                                  %~ requires a complicated initialization 
                                  %~ of the cloned object\\
        %~ \hline
        %~ \highlightbf{Template Method} & large model based on this 
                                        %~ pattern can have 
                                        %~ \highlightit{Factory Method}s 
                                        %~ as single steps\\
        %~ \hline
    %~ \end{tabularx}
%~ \end{table}
    
\newpage
%%% %%% ABSTRACT FACTORY %%% %%% %%%
\subsection{Abstract Factory}
\vspace{-0.5cm}
%%% rephrase better into following table
%~ An Abstract Factory is finally declared with a \textbf{list of creation methods} 
%~ for all the items that are part of the product family. For each variant of a product family, 
%~ \textbf{a separate factory class} based on the AbstractFactory interface will be needed 
%~ and \textbf{clients access with both factories and products via their abstract interfaces}.
\begin{table}[h!t]
    \begin{tabularx}{\textwidth}{|l|X|}
        \hline
        \multirow{1}{*}{\textbf{What for}} & \highlightit{producing families 
                                             of related items} 
                                             without specifying their 
                                             concrete classes\\
        \hline
        \multirow{1}{*}{\textbf{How to}} & declaring 
                                           \highlightit{explicit interfaces} 
                                           for each distinct product 
                                           of the product family,
                                           and making all variants of 
                                           products follow those 
                                           interfaces\\
        \hline
        \textbf{When needed} & $\bullet$ client needs to work with 
                               various families of related objects, \\
                             & \quad without depending on the latter's 
                               concrete classes \\
                             & $\bullet$ a set of \textit{factory methods} 
                               blurring the main responsibility of 
                               a class\\
                             & \quad needs to be managed better \\
        \hline
        \end{tabularx}
\end{table}

\vspace{-0.5cm}
\toclessnparagraph{Structure}
\vspace{-1.5cm}
\begin{table}[h!t]
    \begin{center}
    \begingroup
    \begin{figure}[H]
        \begin{center}
            \scalebox{0.8}
{
    \begin{tikzpicture}[show background grid]
            
        \begin{abstractclass}[text width=6cm]{ABCFactory}{0,0}
            \attributett{...}
            \operationtt{+ createProductA(): ProductA}
            \operationtt{+ createProductB(): ProductB}
        \end{abstractclass}
        \begin{class}[text width=6cm]{Factory1}{0,3}
            \implement{ABCFactory}
            \attributett{...}
            \operationtt{+ createProductA(): ProductA}
            \operationtt{+ createProductB(): ProductB}
        \end{class}
        \begin{class}[text width=6cm]{Factory2}{0,-3.5}
            \implement{ABCFactory}
            \attributett{...}
            \operationtt{+ createProductA(): ProductA}
            \operationtt{+ createProductB(): ProductB}
        \end{class}

        \begin{abstractclass}[text width=3cm]{ABCProductA}{-9.5,-0.5}
        \end{abstractclass}
        \begin{class}[text width=3cm]{ProductA1}{-9.5,1.5}
            \implement{ABCProductA}
        \end{class}
        \begin{class}[text width=3cm]{ProductA2}{-9.5,-3}
            \implement{ABCProductA}
        \end{class}

        \begin{abstractclass}[text width=3cm]{ABCProductB}{-5.5,-0.5}
        \end{abstractclass}
        \begin{class}[text width=3cm]{ProductB1}{-5.5,1.5}
            \implement{ABCProductB}
        \end{class}
        \begin{class}[text width=3cm]{ProductB2}{-5.5,-3}
            \implement{ABCProductB}
        \end{class}

        \begin{class}[text width=4.5cm]{Client}{6.5,-0.3}
            \attributett{- factory: ABCFactory}
            \operationtt{+ Client(f: ABCFactory)}
            \operationtt{+ someOperation()}
        \end{class}
        
        \draw [ umlcd style dashed line , - >] ( Client) -- 
        %~ node [ above , sloped , black ]{$ < <$ instantiate $ > >$} 
        ( ABCFactory) ;
    \end{tikzpicture}
}

        \end{center}
    \end{figure}
    \endgroup
    \vspace{-1cm}
    \begin{tabularx}{\textwidth}{|l|X|}
        \hline
        \highlightbf{ABCProductA}, \highlightbf{ABCProductB} & interfaces for a set or family of related but distinct products\\
        \hline
        \highlightbf{ProductA1}, \highlightbf{ProductB1} & implementations of the concrete variants of abstract products.
                                                                   Each abstract product must be implemented in all given variants.\\
        \hline
        \highlightbf{ABCFactory} & interface for the sets of methods needed to create each of 
                                   the abstract products\\
        \hline
        \highlightbf{Factory1} & implement the creation methods of the abstract factory,
                                                         each creating only the corresponding product variant\\
        \hline
        \highlightbf{Client} & can work with any concrete factory as long as it communicates
                               with the objects through their abstract interfaces.
                               Therefore, the signature of concrete-factories' creation-methods
                               must match the corresponding abstract interfaces'.\\ 
        \hline
    \end{tabularx}
    \caption{Abstract Factory Class Diagram.}
    \label{tab:abs-factory-class-diagram}
    \end{center}
\end{table}

\newpage
\toclessnparagraph{Pros and Cons}
\begin{table}[h!t]
    \begin{tabularx}{\textwidth}{|l|X|}
        \hline
        \textbf{Pros} & $\bullet$\, 
                        \highlightit{interchangeable concrete implementations}
                        without changing the code that \\
                      & \quad uses them\\
                      & $\bullet$\, \highlightit{no tight coupling} 
                        between concrete products and client code\\
                      & $\bullet$\, \highlightit{Single Responsibility Principle} \\
                      & $\bullet$\, \highlightit{Open/Close Principle} \\
        \hline
        \textbf{Cons} & $\bullet$\, unnecessary complexity and work 
                                    in the initial writing \\
                      & $\bullet$\, systems can get more difficult to 
                                    debug and maintain due to higher \\
                      & \quad separation and abstraction \\
        \hline      
    \end{tabularx}
\end{table}
    
%~ \toclessnparagraph{Relation with other Design Patterns}
%~ \begin{table}[h!t]
    %~ \begin{tabularx}{\textwidth}{|l|X|}
        %~ \hline
        %~ \highlightbf{Builder} & instead of returning the product 
                                %~ immediately, \textit{Builder} lets 
                                %~ run some additional construction steps 
                                %~ before releasing its product\\
        %~ \hline
        %~ \highlightbf{Prototype} & can be an alternative option to 
                                  %~ compose \highlightit{Abstract Factory} 
                                  %~ methods and classes\\
        %~ \hline
        %~ \highlightbf{Bridge} & paired with \highlightit{Factory Method} 
                               %~ can help manage scenarios in which 
                               %~ abstraction defined by \highlightit{Bridge} 
                               %~ can only work with specific implementations\\
        %~ \hline
        %~ \highlightbf{Singleton} & can be used to implement 
                                  %~ \highlightit{Abstract Factory}\\
        %~ \hline
    %~ \end{tabularx}
%~ \end{table} 
       
\newpage
%%% %%% BUILDER %%% %%% %%%
\subsection{Builder}
\begin{table}[h!t]
    \begin{tabularx}{\textwidth}{|l|X|}
        \hline
        \multirow{1}{*}{\textbf{What for}} & \highlightit{building complex 
                                             objects step-by-step}, producing 
                                             different types and 
                                             representations using 
                                             the same construction 
                                             steps\\
        \hline
        \multirow{1}{*}{\textbf{How to}} & \highlightit{object's construction code 
                                           is moved into a 
                                           Builder} class, 
                                           that builds the object through 
                                           steps that can be executed 
                                           in series and only if needed. 
                                           Several \texttt{Builder}s 
                                           can be implemented in case 
                                           different ways of building 
                                           make sense. A \highlighttt{Director} 
                                           class can be used to manage 
                                           the order of the building 
                                           steps. \\
        \hline
        \textbf{When needed} & $\bullet$ ctors became telescopic because 
                               of too many optional parameters\\
                             & \quad and need to be removed\\
                             & $\bullet$ code need to be able to create 
                               different representations of some\\
                             & \quad products \\
                             & $\bullet$ a \textit{composite} 
                               object construction needs to be managed 
                               step-by-step\\
        \hline
    \end{tabularx}
\end{table}
\vspace{-0.3cm}
\toclessnparagraph{Structure}
\vspace{-1.5cm}
\begin{table}[h!t]
    \begin{center}
    \begingroup
    \begin{figure}[H]
        \begin{center}
    \begin{tikzpicture}[show background grid]
            
        \begin{abstractclass}[text width=3cm]{ABCBuilder}{0,0}
            \operationtt{+ reset()}
            \operationtt{+ buildStepA()}
            %~ \operationtt{+ buildStepB()}
            \operationtt{+ buildStepZ()}
        \end{abstractclass}
        \begin{class}[text width=4.5cm]{BuilderA}{-3,-4.5}
            \implement{ABCBuilder}
            \attributett{result: ProductA}
            \operationtt{+ reset()}
            \operationtt{+ buildStepA()}
            %~ \operationtt{+ buildStepB()}
            \operationtt{+ buildStepZ()}
            \operationtt{+ getResult(): ProductA}
        \end{class}
        \begin{class}[text width=4.5cm]{BuilderB}{3,-4.5}
            \implement{ABCBuilder}
            \attributett{result: ProductB}
            \operationtt{+ reset()}
            \operationtt{+ buildStepA()}
            %~ \operationtt{+ buildStepB()}
            \operationtt{+ buildStepZ()}
            \operationtt{+ getResult(): ProductB}
        \end{class}


        \begin{class}[text width=4.5cm]{Director}{5,-0.5}
            \attributett{- builder: Builder}
            \operationtt{+ Director(builder)}
            \operationtt{+ changeBuilder(builder)}
            \operationtt{+ make(type)}
        \end{class}

    \end{tikzpicture}
\end{center}
 
    \end{figure}
    \endgroup
    \vspace{-0.5cm}
    \begin{tabularx}{\textwidth}{|l|X|}
        \hline
        \highlightbf{ABCBuilder} & interface that declares the production 
                                   steps common to all the builders\\
        \hline
        \highlightbf{BuilderA}, 
        \highlightbf{BuilderB} & implement the construction steps 
                                 needed for products which do not 
                                 follow the common interface\\
        \hline
        \highlightbf{ProductA}, 
        \highlightbf{ProductB} & resulting objects constructed by 
                                 different builders\\
        \hline
        \highlightbf{Director} & might be useful to manage the calls 
                                 to the steps needed\\
        \hline
        \highlightbf{Client} & can associate a builder with a director 
                               in order to let latter be used for any  
                               further construction. An alternative 
                               way relies on passing the 
                               \highlighttt{Builder} object directly to 
                               the production method.\\
        \hline
    \end{tabularx}
    \caption{Builder Class Diagram.}
    \label{tab:builder-class-diagram}
    \end{center}
\end{table}
    
\newpage
\toclessnparagraph{Pros and Cons}
\begin{table}[h!t]
    \begin{tabularx}{\textwidth}{|l|X|}
        \hline
        \textbf{Pros} & $\bullet$ constrution goes step-by-step and 
                                  can be deferred or run recursively\\
                      & $\bullet$ construction steps can be shared 
                                  in case \\
                      & $\bullet$ \highlightbf{Single Responsibility Principle} \\
        \hline
        \textbf{Cons} & multiple classes required by this pattern 
                        can increase the code complexity \\
        \hline      
    \end{tabularx}
\end{table}

%~ \toclessnparagraph{Relation with other Design Patterns}
%~ \begin{table}[h!t]
    %~ \begin{tabularx}{\textwidth}{|l|X|}
        %~ \hline
        %~ \highlightbf{Abstract Factory} & is similar to \highlightit{Builder} 
                                         %~ but it specializes in 
                                         %~ creating families of related 
                                         %~ objects and return them immediately\\
        %~ \hline
        %~ \highlightbf{Composite} & can be used along with \highlightit{Builder} 
                                  %~ to manage better the construction 
                                  %~ steps\\
        %~ \hline
        %~ \highlightbf{Bridge} & used along with \highlightit{Builder} could 
                               %~ help letting the director class play 
                               %~ the role of the abstraction and different 
                               %~ builders act as implementations\\
        %~ \hline
        %~ \highlightbf{Singleton} & can help implementing a class based on
                                  %~ \highlightit{Builder} that need only an 
                                  %~ instance of it to be created\\ 
        %~ \hline
    %~ \end{tabularx}
%~ \end{table}

\newpage
%%% %%% PROTOTYPE %%% %%% %%%
\subsection{Prototype}
\begin{table}[h!t]
    \begin{tabularx}{\textwidth}{|l|X|}
        \hline
        \multirow{1}{*}{\textbf{What for}} & \highlightit{copying existing objects} 
                                             without implementing code 
                                             which depends on their 
                                             classes\\
        \hline
        \multirow{1}{*}{\textbf{How to}} & a common interface is implemented 
                                           for all objects that need 
                                           to be cloned, and the actual 
                                           cloning procedure is delegated 
                                           to the objects themselves 
                                           in order to have access to 
                                           private and protected details 
                                           as well \\
        \hline
        \textbf{When needed} & $\bullet$ code shouldn't depend on the 
                                         concrete classes of the objects 
                                         being \\
                             & \quad copied\\
                             & $\bullet$ reducing the number of subclasses 
                                         only differing in the way of \\
                             & \quad initializing their respective objects.\\
        \hline
    \end{tabularx}
\end{table}

\toclessnparagraph{Structure}
\vspace{-1.5cm}
\begin{table}[h!t]
    \begin{center}
    \begingroup
    \begin{figure}[H]
        \begin{center}
    \begin{tikzpicture}[show background grid]
        
        \begin{abstractclass}[text width=4.5cm]{ABCPrototype}{0,0}
            \operationtt{+ clone(): Prototype}
        \end{abstractclass}
        \begin{class}[text width=5.5cm]{Prototype}{0,-3}
            \implement{ABCPrototype}
            \attributett{- field1}
            \operationtt{+ ConcretePrototype(prototype)}
            \operationtt{+ clone(): Prototype}

        \end{class}
        \begin{class}[text width=5cm]{SubPrototype}{0,-6}
            \inherit{Prototype}
            \attributett{- field2}
            \operationtt{+ SubPrototype(prototype)}
            \operationtt{+ clone(): Prototype}
        \end{class}
        
        \begin{class}[text width=2cm]{Client}{-5,-0.3}
                    
        \end{class}
    \end{tikzpicture}
\end{center}
        
    \end{figure}
    \endgroup
    \vspace{-0.5cm}
    \begin{tabularx}{\textwidth}{|l|X|}
        \hline
        \highlightbf{ABCPrototype} & interface that declares the 
                                     cloning method\\
        \hline
        \highlightbf{Prototype} & class that implements the cloning 
                                  method; it can have some logic needed 
                                  to manage edge cases\\            
        \hline
        \highlightbf{SubPrototype} & a class that can be handy to add 
                                     subsets of prototypes\\
        \hline
        \highlightbf{Client} & can produce a copy of any object which 
                               follows the prototype interface\\
        \hline
    \end{tabularx}
    \caption{Prototype Class Diagram.}
    \label{tab:prototype-class-diagram}
    \end{center}
\end{table}

\newpage
\toclessnparagraph{Pros and Cons}
\begin{table}[h!t]
    \begin{tabularx}{\textwidth}{|l|X|}
        \hline
        \textbf{Pros} & $\bullet$ cloning objects without coupling 
                                  to their concrete classes \\
                      & $\bullet$ repeated initialization can get 
                                  removed in favor of cloning pre-built\\
                      & \quad prototypes \\
                      & $\bullet$ easier production of complex objects\\
                      & $\bullet$ alternative to inheritance when 
                                  configuring complex objects\\
        \hline
        \textbf{Cons} & cloning can get more difficult when circular 
                        references are present \\
        \hline      
    \end{tabularx}
\end{table}

%~ \toclessnparagraph{Relation with other Design Patterns}
%~ \begin{table}[h!t]
    %~ \begin{tabularx}{\textwidth}{|l|X|}
        %~ \hline
        %~ \highlightbf{Abstract Factory} & can also use \highlightit{Prototype} 
                                         %~ in place of \highlightit{Factory Method} 
                                         %~ to compose its methods\\
        %~ \hline
        %~ \highlightbf{Command} & can work along with \highlightit{Prototype} 
                                %~ to save copies of commands in the history\\
        %~ \hline
        %~ \highlightbf{Composite}, 
        %~ \highlightbf{Decorator} & used along with \highlightit{Prototype}, 
                                  %~ let cloning complex structure instead 
                                  %~ of rebuilding from scratch  \\
        %~ \hline
        %~ \highlightbf{Factory Method} & does not have inheritance 
                                       %~ drawbacks, since it is not 
                                       %~ based on that, but requires a 
                                       %~ complicated initialization\\
        %~ \hline
        %~ \highlightbf{Memento} & can be replaced by \highlightit{Prototype} 
                                %~ if the object's state to store does 
                                %~ not have links to external resources
                                %~ \\
        %~ \hline
        %~ \highlightbf{Singleton} & can help implement a prototype in 
                                  %~ case at most a single instance can 
                                  %~ be present at the same time \\
        %~ \hline
    %~ \end{tabularx}
%~ \end{table} 
    
\newpage
%%% %%% SINGLETON %%% %%% %%%
\subsection{Singleton}
\begin{table}[h!t]
    \begin{tabularx}{\textwidth}{|l|X|}
        \hline
        \multirow{1}{*}{\textbf{What for}} & \highlightit{ensuring that a class could have one instance only}, \\
                                           & while providing a global access point to the instance \\
        \hline
        \multirow{1}{*}{\textbf{How to}} & the \highlightit{ctor is made private} 
                                           to prevent calls from outside the class \\
                                         & and a \highlightit{static creation method} is used in its place in order to create \\
                                         & the object and return it for successive calls.\\
        \hline
        \textbf{When needed} & $\bullet$ a class in a program should have just a single instance available, \\
                             & \quad as it can be needed for database object shared by different parts of \\
                             & \quad the code \\
                             & $\bullet$ a stricter control over global variables is needed \\
        \hline
    \end{tabularx}
\end{table}

\toclessnparagraph{Structure}
\begin{table}[h!t]
    \begin{center}
    \begingroup
    \begin{figure}[H]
    \begin{center}
    \begin{tikzpicture}%[show background grid]

        \begin{class}[text width=5cm]{Singleton}{0,0}
            \attributett{- instance: Singleton}
            \operationtt{- Singleton()}
            \operationtt{- getInstance(): Singleton}
        \end{class}
        
        \begin{class}[text width=3cm]{Client}{-5.5,-.3}
                    
        \end{class}
        
    \end{tikzpicture}
\end{center}
    
    \end{figure}
    \endgroup
    \vspace{-0.5cm}
    \begin{tabularx}{\textwidth}{|l|X|}
        \hline
        \highlightbf{Singleton} & class declaring a \texttt{static} method that return the same instance it creates the 
        first time it's called; this class' ctor is made \texttt{private} to hide it from external calls.\\
        \hline
    \end{tabularx}
    \caption{Singleton Class Diagram.}
    \label{tab:singleton-class-diagram}
    \end{center}  
\end{table}

\toclessnparagraph{Pros and Cons}
\begin{table}[h!t]
    \begin{tabularx}{\textwidth}{|l|X|}
        \hline
        \textbf{Pros} & $\bullet$ avoiding \highlightit{global namespace pollution} 
                        introducing global variables \\
                      & $\bullet$ \highlightit{initialization} takes place 
                        only the first time the object is requested\\
                      & $\bullet$ \highlightit{multiple but limited instances} 
                        can be created, modifing only the static\\
                      & \quad method\\
        \hline
        \textbf{Cons} & $\bullet$ violation of \highlightit{Single Responsibility Principle} \\
                      & $\bullet$ can mask bad design decisions, like too much overlapping between\\
                      & \quad components \\
                      & $\bullet$ multithreaded programs need special care to avoid multiple instances\\
                      & \quad creation \\
                      & $\bullet$ unit tests could be difficult due to the impossibility for mock objects to \\
                      & \quad access the \texttt{Singleton}'s ctor. \\
        \hline      
    \end{tabularx}
\end{table}

%~ \toclessnparagraph{Relation with other Design Patterns}
%~ \begin{table}[h!t]
    %~ \begin{tabularx}{\textwidth}{|l|X|}
        %~ \hline
        %~ \highlightbf{Facade} & can be implemented as a \highlightit{Singleton} 
                               %~ since a single facade object is 
                               %~ usually enough\\
        %~ \hline
        %~ \highlightbf{Flyweight} & resembles \highlightit{Singleton} 
                                  %~ if all the shared states of the objects 
                                  %~ are reduced to one flyweight. 
                                  %~ The \textbf{difference} is: there is 
                                  %~ only one instance for singletons, 
                                  %~ whereas there are almost always 
                                  %~ several flyweights objects; singletons 
                                  %~ are mutable, flyweights are not\\
        %~ \hline
        %~ \highlightbf{Abstract Factories} & can be implemented 
                                           %~ as \highlightit{Singleton}\\
        %~ \highlightbf{Builders}           & \\
        %~ \highlightbf{Prototypes} & \\
        %~ \hline
    %~ \end{tabularx}
%~ \end{table}

\newpage
%%%
%%% BEHAVIORAL PATTERNS %%%
%%%
\section{Behavioral Patterns}
%%% %%% CHAIN OF RESPONSIBILITY %%% %%% %%%
\subsection{Chain of Responsibility}
\begin{table}[h!t]
    \begin{tabularx}{\textwidth}{|l|X|}
        \hline
        \multirow{1}{*}{\textbf{What for}} & \highlightit{pass a request along a 
                                             chain of handlers}, each 
                                             of which can either\\
                                           & manage the request or pass 
                                             it to the next handler 
                                             in the chain\\
        \hline
        \multirow{1}{*}{\textbf{How to}} & each particular behavior 
                                           is implemented into a 
                                           stand-alone handler, \\
                                         & the handlers are linked 
                                           one to another in a chain 
                                           and the request,\\
                                         & along with the data needed, 
                                           is passed along the chain\\
        \hline
        \textbf{When needed} & $\bullet$ different kind of requests 
                               are expected to be processed in 
                               different \\
                             & \quad ways and the type of the request 
                               is not known beforehand \\
                             & $\bullet$ the set or the order of the 
                               handlers can change at run time\\
                             %~ & $\bullet$ \\
                             
        \hline
    \end{tabularx}
\end{table}

\toclessnparagraph{Structure}
\vspace{-1.5cm}
\begin{table}[h!t]
    \begin{center}
    \begingroup
    \begin{figure}[H]
    \begin{center}
\tikzfig{chain_of_responsibility_design_pattern}
\begin{tikzpicture}%[show background grid]
    \begin{abstractclass}[text width=4cm]{ABCHandler}{0,0}
        \operationtt{+setNext(h: Handler)}
        \operationtt{+handle(request)}
        \end{abstractclass}
        
        \begin{class}[text width=4cm]{Handler}{0,-4}
             \implement{ABCHandler}
             \attributett{- next: Handler}      
            \operationtt{+setNext(h: Handler)}
            \operationtt{+handle(request)}
        \end{class}
        \begin{class}[text width=4cm]{SubHandlerA}{-3,-7.5}
             \inherit{Handler}
             \attributett{- next: Handler}      
            \operationtt{+setNext(h: Handler)}
            \operationtt{+handle(request)}
        \end{class}    
        \begin{class}[text width=4cm]{SubHandlerB}{3,-7.5}
             \inherit{Handler}
             \attributett{- next: Handler}      
            \operationtt{+setNext(h: Handler)}
            \operationtt{+handle(request)}
        \end{class}
        
        \begin{class}[text width=3cm]{Client}{-6,-0.5}
        \end{class}
        
        \draw [ - >] ( Client.east) ++ (0,-.25) -- ( ABCHandler.west) ;                  

        
\end{tikzpicture}
\end{center}
       
    \end{figure}
    \endgroup
    \vspace{-0.5cm}
    \begin{tabularx}{\textwidth}{|l|X|}
        \hline
        \highlightbf{ABCHandler} & common interface for all the concrete 
                                   handlers, usually in terms of\\
                                 & a first method for handling the 
                                   requests and second one to set the \\
                                 & next handler in the chain\\
        \hline
        \highlightbf{Handler} & optional base class to manage the 
                                common part of the requests \\
        \hline
        \highlightbf{SubHandlerA}, & classes implementing the actual 
                                     code to process the request\\
        \highlightbf{SubHandlerB} & or pass the request along the 
                                    chain if not possible\\
        \hline
        \highlightbf{Client} & responsible of creating chains, both 
                               statically and dinamically, and\\
                             & of sending the request to the first 
                               or any other handler in the chain\\
        \hline
    \end{tabularx}
    \caption{Chain of Responsibility Class Diagram.}
    \label{tab:chain-of-resp-class-diagram}     
    \end{center}       
\end{table}

\toclessnparagraph{Pros and Cons}
\begin{table}[h!t]
    \begin{tabularx}{\textwidth}{|l|X|}
        \hline
        \textbf{Pros} & $\bullet$ the order of request handling can 
                        be under complete control \\
                      & $\bullet$ flexibility in deciding whether a 
                        handler should stop the request or pass it \\
                      & \quad along the chain \\
                      & $\bullet$ \highlightit{Single Responsibility Principle} \\
                      & $\bullet$ \highlightit{Open/Close Principle}\\
        \hline
        \textbf{Cons} & some requests may end up unhandled \\
        \hline      
    \end{tabularx}
\end{table}
    
%~ \toclessnparagraph{Relation with other Design Patterns}
%~ \begin{table}[h!t]
    %~ \begin{tabularx}{\textwidth}{|l|X|}
        %~ \hline
        %~ \highlightbf{Chain of Responsibility} & passes a request along a dynamic chain of potential\\
                                              %~ & receiver;\\
        %~ \highlightbf{Command}                 & makes unidirectional connections between senders and receivers;\\
        %~ \highlightbf{Mediator}                & forces senders and receivers to communicate indirectly through a mediator;\\
        %~ \highlightbf{Observer}                & lets receivers un/subscribe from a service. \\
        %~ \hline
        %~ \highlightbf{Composite} & when used along with 
                                  %~ \highlightit{CoR} 
                                  %~ can help building a tree of responsibilities\\
        %~ \hline
        %~ \highlightbf{Command} & can help implement the chain handler, 
                                %~ letting the user executing the same 
                                %~ operation in a series of different contextes\\
        %~ \hline
        %~ \highlightbf{Decorator} & has a similar class structure, relying 
                                  %~ on recursive composition through a 
                                  %~ series of objects; \highlightit{CoR} 
                                  %~ lets execute operations independently 
                                  %~ of each other, whereas 
                                  %~ \highlightit{Decorator} extends an 
                                  %~ object behavior keeping it consistent 
                                  %~ with a base interface\\
        %~ \hline
    %~ \end{tabularx}
    % \caption{Chain of Responsibility relation with other Patterns.}
    %~ \hypertarget{tab:cor-relation-with-others}{}  
%~ \end{table}
    
\newpage
%%% %%% COMMAND %%% %%% %%%
\subsection{Command}
\begin{table}[h!t]
    \begin{tabularx}{\textwidth}{|l|X|}
        \hline
        \multirow{1}{*}{\textbf{What for}} & \highlightit{turning a request into a 
                                             stand-alone object} containing 
                                             all information about it.\\
        \hline
        \multirow{1}{*}{\textbf{How to}} & applying the \highlightit{Principle of Separation of Concerns}\\
        \hline
        \textbf{When needed} & $\bullet$ there is the need to parametrize 
                               objects with operations\\
                             & $\bullet$ there is the need for 
                               reversible operations \\
        \hline
    \end{tabularx}
\end{table}

\toclessnparagraph{Structure}
\vspace{-1.5cm}
\begin{table}[h!t]
    \begin{center}
    \begingroup
    \begin{figure}[H]
    \begin{center}
    \begin{tikzpicture}[show background grid]
        
        \begin{class}[text width=4cm]{Invoker}{0,0}
            \attributett{- command}
            \operationtt{+ setCommand(command)}
            \operationtt{+ executeCommand()}
        \end{class}
        
        \begin{abstractclass}[text width=3cm]{ABCCommand}{5,-0.1}
            \operationtt{+ execute()}
        \end{abstractclass}
        \begin{class}[text width=5.5cm]{CommandA}{2.3,-3}
            \implement{ABCCommand}
            \attributett{- receiver}
            \attributett{- params}
            \attributett{+ CommandA(receiver, params)}
            \attributett{+ execute()}
        \end{class}
        \begin{class}[text width=3cm]{CommandB}{7.7,-3}
            \implement{ABCCommand}
            \operationtt{+ execute()}
        \end{class}
        
        
        
        \begin{class}[text width=3cm]{Client}{-5,-.5}
                    
        \end{class}
        \begin{class}[text width=4cm]{Receiver}{-5,-4}
            \attributett{...}
            \operationtt{+ operation(a, b, c)}
        \end{class}
        
        \draw [ umlcd style , - >] ( Client)  -- ( Receiver);
        \draw [ umlcd style , - >] ( Invoker)  -- ( ABCCommand);
        \draw [ umlcd style dashed line, - >] ( Client)  -- ( CommandA);
    \end{tikzpicture}
\end{center}
   
    \end{figure}
    \endgroup
\vspace{-0.5cm}
    \begin{tabularx}{\textwidth}{|l|X|}
        \hline
        \highlightbf{Invoker}/\highlightbf{Sender} & class responsible 
                                                     for creating and 
                                                     sending requests, 
                                                     directly triggering 
                                                     a reference 
                                                     to a command 
                                                     object, usually 
                                                     provided and not 
                                                     created by the 
                                                     sender\\
        \hline
        \highlightbf{ABCCommand} & interface declaring the method to 
                                   execute the command\\
        \hline
        \highlightbf{CommandA}, & classes implementing the concrete 
                                 requests, passin the call to one\\
        \highlightbf{CommandB} & of the business logic objects\\
        \hline
        \highlightbf{Receiver} & class that implements the business 
                                 logic and does the actual work 
                                 most of the times\\
        \hline
        \highlightbf{Client} & class responsible for creating and 
                               configuring commands, passing all the 
                               parameters needed\\
        \hline
    \end{tabularx}    
    \caption{Command Class Diagram.}
    \label{tab:command-class-diagram}     
    \end{center}
\end{table}

\newpage
\toclessnparagraph{Pros and Cons}
\begin{table}[h!t]
    \begin{tabularx}{\textwidth}{|l|X|}
        \hline
        \textbf{Pros} & $\bullet$ \highlightit{Single Responsibility Principle}\\
                      & $\bullet$ \highlightit{Open/Close Principle}\\
                      & $\bullet$ implementing \textit{undo/redo} is easier\\
                      & $\bullet$ implementing deferred execution of operation is easier\\
                      & $\bullet$ set of simple commands can be assebled in sets\\
        \hline
        \textbf{Cons} & programs end up having more layers \\
        \hline      
    \end{tabularx}
\end{table}
    
%~ \toclessnparagraph{Relation with other Design Patterns}
%~ \begin{table}[h!t]
    %~ \begin{tabularx}{\textwidth}{|l|X|}
        %~ \hline
        %~ \highlightbf{Chain of Responsibility} & same considerations 
                                                % in~\ref{tab:cor-relation-with-others} \\ 
                                                %~ for~\hyperlink{tab:cor-relation-with-others}{CoR} \\ 
        %~ \highlightbf{Command} & \\
        %~ \highlightbf{Mediator} & \\
        %~ \highlightbf{Observer} & \\
        %~ \hline
        %~ \highlightbf{Chain of Responsibility} & handlers can be 
                                                %~ implemented in terms 
                                                %~ of \highlightit{Command} \\
        %~ \hline
        %~ \highlightbf{Memento} & can be used along with \highlightit{Command} 
                                %~ to implement an undo routines\\
        %~ \hline
        %~ \highlightbf{Strategy} & also can be used to parameterize an 
                                 %~ object with some action. The \textbf{difference} 
                                 %~ is that \highlightit{Command} convert 
                                 %~ any operation into an object whereas 
                                 %~ \highlightit{Strategy} is generally 
                                 %~ used to achieve the same result with 
                                 %~ different algorithms \\
        %~ \hline
        %~ \highlightbf{Prototype} & can help save copies of commands 
                                  %~ into history\\
        %~ \hline
        %~ \highlightbf{Visitor} & can be seen as a powerful version of 
                                %~ \highlightit{Command} because it can 
                                %~ execute operations over various objects 
                                %~ of different classes\\
        %~ \hline
    %~ \end{tabularx}
    %~ \hypertarget{tab:command-relation-with-others}{}
%~ \end{table} 
        
\newpage
%%% %%% ITERATOR %%% %%% %%%
\subsection{Iterator}
\begin{table}[h!t]
    \begin{tabularx}{\textwidth}{|l|X|}
        \hline
        \multirow{1}{*}{\textbf{What for}} & \highlightit{traversing 
                                             elements of a collection 
                                             without exposing its underlying} \\
                                           & \highlightit{representation}\\
        \hline
        \multirow{1}{*}{\textbf{How to}} & encapsulating the details 
                                           of the traversal into a 
                                           object called iterator \\
        \hline
        \textbf{When needed} & $\bullet$ collection with complex 
                               data structure not to expose to users\\
                             & $\bullet$ reducing code duplication 
                               for traversals \\
                             & $\bullet$ traversing data structures, 
                               that can be unknown beforehand\\
        \hline
    \end{tabularx}
\end{table}
\vspace{-0.5cm}
%~ \newpage
\toclessnparagraph{Structure}
\vspace{-1.5cm}
\begin{table}[h!t]
    \begin{center}
    \begingroup
    \begin{figure}[H]
        \begin{center}
    \begin{tikzpicture}%[show background grid]

        \begin{abstractclass}[text width=3.5cm]{ABCIterator}{0,0}
            \operationtt{+ getNext()}
            \operationtt{+ hasMore(): bool}
        \end{abstractclass}
        \begin{class}[text width=4.5cm]{IteratorA}{0,-4}
            \implement{ABCIterator}
            \attributett{- collection: Collection}
            \attributett{- iterationState}
            \operationtt{+ IteratorA(c: Collection)}
            \operationtt{+ getNext()}
            \operationtt{+ hasMore(): bool}
        \end{class}

        \begin{abstractclass}[text width=4.5cm]{ABCCollection}{6.5,0}
            \operationtt{+ createIterator: ABCIterator}
        \end{abstractclass}
        \begin{class}[text width=4.5cm]{ConcreteCollection}{6.5,-4}
            \implement{ABCCollection}
            \attributett{...}
            \operationtt{...}
            \operationtt{+ createIterator(): ABCIterator}
        \end{class}
        
        \begin{class}[text width=3cm]{Client}{3,2}

        \draw [ - >] ( Client) ++ (+8,2) -- ( ABCCollection) ;                                   
        \draw [ - >] ( Client.south) ++ (+8.3,3) -- ( ABCIterator) ;                                   



        \end{class}

    \end{tikzpicture}
\end{center}
       
    \end{figure}
    \endgroup
    \vspace{-0.5cm}
    \begin{tabularx}{\textwidth}{|l|X|}
        \hline
        \highlightbf{ABCIterator} & interface declaring the operations 
                                    needed to traverse a collection, 
                                    like fetching the next element, 
                                    keeping track of the current 
                                    position and so on\\
        \hline
        \highlightbf{IteratorA} & implements specific algorithms, 
                                  independent from each others in 
                                  order to be used in any number on 
                                  the same collection\\
        \hline
        \highlightbf{ABCCollection} & interface declaring one or multiple 
                                      methods for getting iterator 
                                      compatible with the collection. 
                                      The return type of these methods 
                                      must match the Iterator interface.\\
        \hline
        \highlightbf{CollectionA} & returns instances of a particular 
                                    concrete iterator class on 
                                    client's request.\\
        \hline
        \highlightbf{Client} & can work with both the collection and 
                               the iterator using their interfaces.\\
        \hline
    \end{tabularx}
    \caption{Iterator Class Diagram.}
    \label{tab:iterator-class-diagram}     
    \end{center}          
\end{table}  

\newpage
\toclessnparagraph{Pros and Cons}
\begin{table}[h!t]
    \begin{tabularx}{\textwidth}{|l|X|}
        \hline
        \textbf{Pros} & $\bullet$ \highlightit{Single Responsibility Principle} \\
                      & $\bullet$ \highlightit{Open/Close Principle} \\
                      & $\bullet$ parallel iterations are possible\\
                      & $\bullet$ iteration can be delayed to when it is needed\\
        \hline
        \textbf{Cons} & can be an overkill for simple collections \\
                      & can be less efficient than going directly 
                        through the elements of some specialized 
                        collection\\
        \hline      
    \end{tabularx}
\end{table}

%~ \toclessnparagraph{Relation with other Design Patterns}
%~ \begin{table}[h!t]
    %~ \begin{tabularx}{\textwidth}{|l|X|}
        %~ \hline
        %~ \highlightbf{Composite} & elements can be traversed with the 
                                  %~ help of \highlightit{Iterator} \\
        %~ \hline
        %~ \highlightbf{Factory Method} & help collection subclasses return 
                                       %~ different type of iterators\\
        %~ \hline
        %~ \highlightbf{Memento} & can help capture the state before an 
                                %~ iteration step \\
        %~ \hline
        %~ \highlightbf{Visitor} & can help traverse complex data structure 
                                %~ and perform some operation over the 
                                %~ elements\\
        %~ \hline
    %~ \end{tabularx}
%~ \end{table} 

\newpage
%%% %%% MEDIATOR %%% %%% %%%
\subsection{Mediator}
\begin{table}[h!t]
    \begin{tabularx}{\textwidth}{|l|X|}
        \hline
        \multirow{1}{*}{\textbf{What for}} & \highlightit{reducing chaotic and 
                                             confusing dependencies} 
                                             between objects and 
                                             \highlightit{restricting direct 
                                             communications} between 
                                             the objects \\
        \hline
        \multirow{1}{*}{\textbf{How to}} & \\
        \hline
        \textbf{When needed} & $\bullet$ changing an already exhisting 
                               class which is strongly coupled \\
                             & $\bullet$ reusing a strongly dependent 
                               component in a different program \\
                             & $\bullet$ reusing basic behaviours 
                               in completely different contexts\\
        \hline
    \end{tabularx}
\end{table}

%~ \begin{itemize}
    %~ \item \textbf{Components} are classes containing some business logic. 
    %~ Each one has a reference to a mendiator, whose referenced type is 
    %~ the Mediator interface's. Components are unaware of the actual Meditor class, 
    %~ so that it could be reused in other contexts with different mediators.
    %~ \item \textbf{Mediator} interface declares the communication structure 
    %~ between components. Components may pass any context as argument of their methods, 
    %~ although they must avoid coupling the receiving component with the sender.
    %~ \item \textbf{Concrete Mediator} implement relation between various components. 
    %~ It could be useful to keep references to all the components managed 
    %~ and their lifecycle as well. 
%~ \end{itemize}
%~ Components are not aware of each other. In case some communication needs 
%~ to go from a Component to a different one, it must pass through the Mediator. 
%~ Likewise, the sender doesn't know who will handle the request as well as the 
%~ receiver doesn't know who sent the request.

\toclessnparagraph{Structure}
\vspace{-1.5cm}
\begin{table}[h!t]
    \begin{center}
    \begingroup
    \begin{figure}[H]
        \begin{center}
    \begin{tikzpicture}[show background grid]

        \begin{abstractclass}[text width=3.5cm]{ABCMediator}{0,0}
            \operationtt{+ notify(sender)}
        \end{abstractclass}
        \begin{class}[text width=3.5cm]{Mediator}{0,-3}
            \implement{ABCMediator}
            \attributett{- componentA}
            \attributett{- componentA}
            \attributett{- componentA}
            \attributett{- componentA}
            \operationtt{+ notify(sender)}
            \operationtt{+ reactOnA()}
            \operationtt{+ reactOnB()}
            \operationtt{+ reactOnC()}
        \end{class}

        \begin{abstractclass}[text width=4.5cm]{ABCComponent}{8,0}
            \operationtt{+ operation()}
        \end{abstractclass}
        
        \begin{class}[text width=3cm]{Component}{4,-4}
            \implement{ABCComponent}
            \attributett{- m: ABCMediator}       
            \operationtt{+ operation()}
        \end{class}
        \begin{class}[text width=3cm]{Component}{8,-4}
            \implement{ABCComponent}
            \attributett{- m: ABCMediator}       
            \operationtt{+ operation()}                    
        \end{class}
        \begin{class}[text width=3cm]{Component}{12,-4}
            \implement{ABCComponent}
            \attributett{- m: ABCMediator}       
            \operationtt{+ operation()}                    
        \end{class}

    \end{tikzpicture}

\end{center}
       
    \end{figure}
    \endgroup
    \begin{tabularx}{\textwidth}{|l|X|}
        \hline
        \highlightbf{Component} & interface for communication with other Components through its Mediator \\
        \hline
        \highlightbf{ComponentA}, \dots, & implements the Colleague interface and communicates \\
        \highlightbf{ComponentC} & with other Components through its Mediator \\
        \hline
        \highlightbf{Mediator} & interface for communication between Component objects\\
        \hline
        \highlightbf{ConcMediator} & implements the mediator interface, coordinating communication between Component objects. \\
                                   & It is aware of all of the Component and their purposes with regards to inter-communication\\
        \hline
    \end{tabularx}    
    \caption{Mediator Class Diagram.}
    \label{tab:mediator-class-diagram}     
    \end{center} 
\end{table}    

\newpage
\toclessnparagraph{Pros and Cons}
\begin{table}[h!t]
    \begin{tabularx}{\textwidth}{|l|X|}
        \hline
        \textbf{Pros} & \highlightit{Single Responsibility Principle} \\
                      & \highlightit{Open/Close Principle} \\
                      & reducing coupling between various components \\
                      & reusing individual components is easier \\
        \hline
        \textbf{Cons} & mediator can evolve into a
                        \href{https://en.wikipedia.org/wiki/God_object}{God object} \\
        \hline      
    \end{tabularx}
\end{table}

\newpage
%~ \toclessnparagraph{Relation with other Design Patterns}
%~ \begin{table}[h!t]
    %~ \begin{tabularx}{\textwidth}{|l|X|}
        %~ \hline
        %~ \highlightbf{Chain of Responsibility} & same considerations 
                                                %~ for~\hyperlink{tab:cor-relation-with-others}{CoR} \\ 
        %~ \highlightbf{Command} & \\
        %~ \highlightbf{Mediator} & \\
        %~ \highlightbf{Observer} & \\
        %~ \hline
        %~ \highlightbf{Facade} & helps organizing collaboration between 
                               %~ tightly coupled classes, achieving it 
                               %~ without introducing any new functionality 
                               %~ and letting objects in the subsystem 
                               %~ communicate directly. \\ 
                             %~ & \highlightit{Mediator} centralizes 
                               %~ the communications, avoiding the objects 
                               %~ communicate only directly\\
        %~ \hline
        %~ \highlightbf{Observer} & estabilishes dynamic one-way connections 
                                 %~ between objects whereas \highlightbf{Mediator} 
                                 %~ aliminates mutual dependencies among 
                                 %~ a set of system components. \\
                               %~ & One implementation of \highlightit{Mediator} 
                                 %~ relies on \highlightit{Observer}: 
                                 %~ the former is the publisher and 
                                 %~ the components act as subscribers 
                                 %~ to mediator's events.\\
        %~ \hline
    %~ \end{tabularx}
    % \caption{Mediator relation with other Patterns.}
    % \caption{Mediator relation with other Patterns.}
    %~ \hypertarget{tab:mediator-relation-with-others}{}
%~ \end{table} 

\newpage
%%% %%% MEMENTO %%% %%% %%%
\subsection{Memento}
\begin{table}[h!t]
    \begin{tabularx}{\textwidth}{|l|X|}
        \hline
        \multirow{1}{*}{\textbf{What for}} & \highlightit{saving/restoring 
                                             states} of an object 
                                             without exposing 
                                             internals\\
        \hline
        \multirow{1}{*}{\textbf{How to}} & \\
        \hline
        \textbf{When needed} & $\bullet$ snapshots of an object's 
                               state are needed in order to restore 
                               them later \\
                             & $\bullet$ when encapsulation would be 
                               violated by directly accessing object's 
                               fields\\
        \hline
    \end{tabularx}
\end{table}

%~ In the previous design, the Memento class is embedded in the Originator, 
%~ in order to make it easier to the latter to access the former's methods and data members.

\toclessnparagraph{Structure}
%~ \vspace{-1.5cm}
\begin{table}[h!t]
    \begin{center}
    \begingroup
    \begin{figure}[H]
        \begin{center}
    \begin{tikzpicture}%[show background grid]

        \begin{class}[text width=4cm]{Originator}{0,0}
            \attributett{- state}
            \operationtt{+ save(): Memento}
            \operationtt{+ restore(m : Memento)}
        \end{class}

        \begin{class}[text width=3cm]{Memento}{5.5,0}
            \attributett{- state}
            \operationtt{- Memento(state)}
            \operationtt{- getState()}
        \end{class}
        
        \begin{class}[text width=4cm]{Caretaker}{11,0}
            \attributett{- originator}
            \attributett{-history: Memento}
            \operationtt{+ doSomething()}
            \operationtt{+ undo()}
        \end{class}

        \draw [umlcd style dashed line , - >] ( Originator.east) -- ( Memento);
        \aggregation{Caretaker}{}{}{Memento}
        
        
    \end{tikzpicture}
    

\end{center}

    \end{figure}
    \endgroup
    \begin{tabularx}{\textwidth}{|l|X|}
        \hline
        \highlightbf{Originator} & generic class that can produce 
                                   snapshots of its own state, and 
                                   restore its state from one of them\\
        \hline
        \highlightbf{Memento} & class that acts like a shapshot, and 
                                could be made immutable\\
        \hline
        \highlightbf{Caretaker} & knows the \textit{why} and \textit{when} 
                                  to capture a snapshot or restore a 
                                  previous one. This class can manage 
                                  an history of Originator's' states 
                                  in a stack data structure.\\
        \hline
    \end{tabularx}
    \caption{Memento Class Diagram.}
    \label{tab:memento-class-diagram}     
    \end{center}     
\end{table}
 
\toclessnparagraph{Pros and Cons}
\begin{table}[h!t]
    \begin{tabularx}{\textwidth}{|l|X|}
        \hline
        \textbf{Pros} & $\bullet$ producing snapshots of the object's state \\
        \hline
        \textbf{Cons} & $\bullet$ memory compsuntion increases \\
                      & $\bullet$ to destroy outdated mementos, the 
                        originator's lifecycle should be tracked\\
                      & $\bullet$ modern dynamic languages do not 
                        guarantee the state within the memento \\
                      & \quad stays untouched\\
        \hline      
    \end{tabularx}
\end{table}

%~ \newpage
%~ \toclessnparagraph{Relation with other Design Patterns}
%~ \begin{table}[h!t]
    %~ \begin{tabularx}{\textwidth}{|l|X|}
        %~ \hline
        %~ \highlightbf{Command} & makes it easier to implement \textit{undo}\\
        %~ \hline
        %~ \highlightbf{Iterator} & can help capture the current iteration 
                                 %~ state and roll it back if necessary\\
        %~ \hline
        %~ \highlightbf{Prototype} & can be a simple alternative if the 
        %~ object state does not have links to external resources\\
        %~ \hline
    %~ \end{tabularx}    
%~ \end{table} 

\newpage
%%% %%% OBSERVER %%% %%% %%%
\subsection{Observer}
\begin{table}[h!t]
    \begin{tabularx}{\textwidth}{|l|X|}
        \hline
        \multirow{1}{*}{\textbf{What for}} & implementing a 
                                             \highlightit{subscription mechanism}\\
        \hline
        \multirow{1}{*}{\textbf{How to}} & \\
        \hline
        \textbf{When needed} & $\bullet$ changes to the state of one 
                               object may require changing other objects 
                               that are unknown beforehand or change 
                               dynamically\\
                             & $\bullet$ objects must observe each 
                               others, for a limited time and in 
                               specific circumstances\\
        \hline
    \end{tabularx}
\end{table}

\toclessnparagraph{Structure}
\vspace{-1cm}
\begin{table}[h!t]
    \begin{center}
    \begingroup
    \begin{figure}[H]
    \begin{center}

    \begin{tikzpicture}[show background grid]

        \begin{class}[text width=6cm]{Publisher}{0,0}
            \attributett{- subscriber: ABCSubscriber}
            \attributett{- mainState}
            \operationtt{+ subscribe(s: ABCSubscriber)}
            \operationtt{+ unsubscribe(s: ABCSubscriber)}
            \operationtt{+ notifySubscriber()}
            \operationtt{+ mainBusinessLogic()}
        \end{class}
        
        \begin{abstractclass}[text width=3.5cm]{ABCSubscriber}{7.5,0}
            \operationtt{+ update(context)}
        \end{abstractclass}
        \begin{class}[text width=3.5cm]{SubscriberA}{5.5,-3}
            \implement{ABCSubscriber}
            \attributett{...}
            \operationtt{+ update(context)}
        \end{class}
        \begin{class}[text width=3.5cm]{SubscriberB}{9.5,-3}
            \implement{ABCSubscriber}
            \attributett{...}
            \operationtt{+ update(context)}
        \end{class}
        
        \begin{class}[text width=3cm]{Client}{0,-5.5}
        \end{class}

    \end{tikzpicture}

\end{center}

\end{figure}
\endgroup
    \begin{tabularx}{\textwidth}{|l|X|}
        \hline
        \highlightbf{Publisher} & class that issues events of interests 
                                  to other objects. This class contains a subscription infrastructure that 
                                  lets subscribers both join and leave the subscriber list \\
        \hline
        \highlightbf{Subscriber} & interface declares the notification interface, that 
                                   in most cases is a single update method with some contextual parameters.\\
        \hline
        \highlightbf{SubscriberA}, \highlightbf{SubscriberB} & implement the different logics in response 
                                                               to notifications issued by the publisher.\\
        \hline
        \highlightbf{Client} & creates both publisher and subscriber objects and 
                               manage the registrations for updates\\
        \hline
    \end{tabularx}
    \caption{Observer Class Diagram.}
    \label{tab:observer-class-diagram}     
    \end{center}     
\end{table}

\newpage
\toclessnparagraph{Pros and Cons}
\begin{table}[h!t]
    \begin{tabularx}{\textwidth}{|l|X|}
        \hline
        \textbf{Pros} & $\bullet$ \highlightbf{Open/Close Principle} \\
        & $\bullet$ relations between objects can be established at runtime \\
        \hline
        \textbf{Cons} & $\bullet$ subscribers can be notified in random order \\
        \hline      
    \end{tabularx}
\end{table}

%~ \newpage
%~ \toclessnparagraph{Relation with other Design Patterns}
%~ \begin{table}[h!t]
    %~ \begin{tabularx}{\textwidth}{|l|X|}
        %~ \hline
        %~ \highlightbf{Chain of Responsibility} & same considerations 
                                                %~ for~\hyperlink{tab:cor-relation-with-others}{CoR} \\ 
        %~ \highlightbf{Command} & \\
        %~ \highlightbf{Mediator} & \\
        %~ \highlightbf{Observer} & \\
        %~ \hline
        %~ \highlightbf{Mediator} & same considerations for~\hyperlink{tab:mediator-relation-with-others}{Mediator} \\
        %~ \hline   
    %~ \end{tabularx}
%~ \end{table}

\newpage
%%% %%% STATE %%% %%% %%%
\subsection{State}
\begin{table}[h!t]
    \begin{tabularx}{\textwidth}{|l|X|}
        \hline
        \multirow{1}{*}{\textbf{What for}} & implementing 
                                             \highlightit{objects whose 
                                             behavior depends on their 
                                             internal state}\\
        \hline
        \multirow{1}{*}{\textbf{How to}} & \\
        \hline
        \textbf{When needed} & $\bullet$ implementing an object that 
                                         behaves differently depending 
                                         on its current state, 
                                         the number of states is 
                                         considerable, and the 
                                         state-specific code changes 
                                         frequently\\
                             & $\bullet$ managing better a class that 
                                         changes its behavior depending 
                                         on some interal field actual 
                                         value\\
                             & $\bullet$ managing better duplicate code 
                                         across similar states and 
                                         transition od a condition-besed 
                                         state machine\\
        \hline
    \end{tabularx}
\end{table}
    
\toclessnparagraph{Structure}
\vspace{-1.5cm}
\begin{table}[h!t]
    \begin{center}
    \begingroup
    \begin{figure}[H]
        \begin{center}
    \begin{tikzpicture}[show background grid]
        
        \begin{class}[text width=4.5cm]{Context}{0,0}
            \attributett{- state}
            \operationtt{+ Context(initialState)}
            \operationtt{+ changeState(state)}
            \operationtt{+ doThis()}
            \operationtt{+ doThat()}
        \end{class}
        
        \begin{abstractclass}[text width=3cm]{ABCState}{8,0}
            \operationtt{+ doThis()}
            \operationtt{+ doThat()}
        \end{abstractclass}
        \begin{class}[text width=4cm]{StateA}{5,-4}
            \implement{ABCState}
            \attributett{- context}
            \operationtt{+ setContext(context)}
            \operationtt{+ doThis()}
            \operationtt{+ doThat()}
        \end{class}
        \begin{class}[text width=4cm]{StateB}{11,-4}
            \implement{ABCState}
            \attributett{- context}
            \operationtt{+ setContext(context)}
            \operationtt{+ doThis()}
            \operationtt{+ doThat()}
        \end{class}
        
        \begin{class}[text width=3cm]{Client}{0,-4.5}
        \end{class}
        
    \end{tikzpicture}
\end{center}
        
    \end{figure}
    \endgroup
    \begin{tabularx}{\textwidth}{|l|X|}
        \hline
        \highlightbf{Context} & stores a reference to one of the concrete state objects, 
                                to which it delegates all the state-related work. The Context interacts with the 
                                state interface, in order to be able to interact with any Concrete State\\
        \hline
        \highlightbf{State} & interface declares the state generic methods, 
                              that should make sense for all concrete states for a matter of consistency\\
        \hline
        \highlightbf{StateA} & implement their own specific methods, 
                               and can make use of abstract  \\
        \highlightbf{StateB} & classes to encapsulate some common behavior\\
        \hline
    \end{tabularx}
    %~ \caption{Observer Class Diagram.}
    \label{tab:observer-class-diagram}     
    \end{center}
\end{table}    

\newpage
\toclessnparagraph{Pros and Cons}
\begin{table}[h!t]
    \begin{tabularx}{\textwidth}{|l|X|}
        \hline
        \textbf{Pros} & $\bullet$ \highlightit{Single Responsibility Principle} \\
                      & $\bullet$ \highlightit{Open/Close Principle} \\
                      & $\bullet$ bulky state machine code is eliminated \\
        \hline
        \textbf{Cons} & can be an overkill in case the state machine 
                        has a few states that change sporadically\\
        \hline     
    \end{tabularx}
\end{table}
    
%~ \toclessnparagraph{Relation with other Design Patterns}
%~ \begin{table}[h!t]
    %~ \begin{tabularx}{\textwidth}{|l|X|}
        %~ \hline
        %~ \highlightbf{Bridge} & have similar structures based on 
                               %~ composition and delegation of the work 
                               %~ \\ 
        %~ \highlightbf{State} & to other objects. \\
        %~ \highlightbf{Strategy} & They solve different 
                                 %~ problems though.\\
        %~ \hline
        %~ \highlightbf{Strategy} & similarly to \highlightit{State} is 
                                 %~ based on composition and changes the 
                                 %~ behavior of the context by delegating 
                                 %~ some work to helper objects. 
                                 %~ The difference is \\
        %~ \hline  
    %~ \end{tabularx}    
    %~ \hypertarget{tab:state-relation-with-others}{}
%~ \end{table} 

\newpage
%%% %%% STRATEGY %%% %%% %%%
\subsection{Strategy}
\begin{table}[h!t]
    \begin{tabularx}{\textwidth}{|l|X|}
        \hline
        \multirow{1}{*}{\textbf{What for}} & defining a \highlightit{family of 
                                             algorithms, each 
                                             implemented in a seprate 
                                             class}, in oder to use 
                                             them interchangeably \\
        \hline
        \multirow{1}{*}{\textbf{How to}} & \\
        \hline
        \textbf{When needed} & $\bullet$ different variants of the same 
                                         algorithm must be used by an 
                                         external object that can be 
                                         able to switch between then \\
                             & $\bullet$ managing better lot of similar 
                                         classes only differing in 
                                         the way they execute some 
                                         behavior \\
                             & $\bullet$ isolating business logic 
                                         from the implementation 
                                         details \\
                             & $\bullet$ better management of massive 
                                         conditional statements that 
                                         switches between variants of 
                                         the same algorithm \\
        \hline
    \end{tabularx}
\end{table}
    
%~ \begin{itemize}
    %~ \item \textbf{Context} has a reference to the Strategy Interface, 
    %~ used to reference the Concrete Strategy in use, and methods to 
    %~ replace the strategy at run-time.
    %~ \item \textbf{Concrete Strategies} 
    %~ \item \textbf{Client} manages the creation of the strategies and pass them 
    %~ to the context.
%~ \end{itemize}

\toclessnparagraph{Structure}
\vspace{-1.5cm}
\begin{table}[h!t]
    \begin{center}
    \begingroup
    \begin{figure}[H]
        \begin{center}
    
    \begin{tikzpicture}[show background grid]
        
        \begin{class}[text width=4.5cm]{Context}{0,0}
            \attributett{- strategy}
            \operationtt{- setStrategy(strategy)}
            \operationtt{- doSomething()}
        \end{class}
        
        \begin{abstractclass}[text width=3cm]{ABCStrategy}{7,0}
            \operationtt{+ execute(data)}
        \end{abstractclass}
        \begin{class}[text width=3cm]{StrategyA}{5,-3}
            \implement{ABCStrategy}
            \operationtt{+ execute(data)}
        \end{class}
        \begin{class}[text width=3cm]{StrategyB}{9,-3}
            \implement{ABCStrategy}
            \operationtt{+ execute(data)}
        \end{class}
        
        \begin{class}[text width=3cm]{Client}{0,-3.2}
        \end{class}
    
    \end{tikzpicture}
    
\end{center}
        
    \end{figure}
    \endgroup
    \begin{tabularx}{\textwidth}{|l|X|}
        \hline
        \highlightbf{Context} & has a reference to the Strategy Interface, 
                                used to reference the Concrete Strategy in use, and methods to 
                                replace the strategy at run-time.\\
        
        \hline
        \highlightbf{Strategy} & interface define the part which is common 
                                 to all concrete strategies and declares 
                                 a method used by the Context 
                                 to execute it.\\
        \hline
        \highlightbf{StrategyA}, \highlightbf{StrategyB} & implement a variety of algorithms 
    needed by the context.\\
        \hline
        \highlightbf{Client} & manages the creation of the strategies and pass them 
                               to the context.\\
        \hline
    \end{tabularx}    
    \caption{Strategy Class Diagram.}
    \label{tab:strategy-class-diagram}     
    \end{center}
\end{table}  
 
\newpage 
\toclessnparagraph{Pros and Cons}
\begin{table}[h!t]
    \begin{tabularx}{\textwidth}{|l|X|}
        \hline
        \textbf{Pros} & $\bullet$ swapping algorithms at runtime \\
                      & $\bullet$ implementation details of an 
                                  algorithm is separate from the client 
                                  using it\\
                      & $\bullet$ inheritance is replaced with composition\\
                      & $\bullet$ \highlightbf{Open/Close Principle}\\
        \hline
        \textbf{Cons} & $\bullet$ can be an overkill in case not many 
                                  algorithms will be used and they will 
                                  not change much at runtime \\
                      & $\bullet$ clients must know the differences 
                                  between the algorithms to be able 
                                  to choose \\
                      & $\bullet$ most of the modern programming 
                                  languages implement different 
                                  versions of an algorithm inside 
                                  a set of anonymous functions and 
                                  let the client use them in a 
                                  similar way\\
        \hline      
    \end{tabularx}
\end{table}

%~ \toclessnparagraph{Relation with other Design Patterns}
%~ \begin{table}[h!t]
    %~ \begin{tabularx}{\textwidth}{|l|X|}
        %~ \hline
        %~ \highlightbf{Bridge} & same consideration 
                               %~ for~\hyperlink{tab:state-relation-with-others}{State} \\
        %~ \highlightbf{State} &  \\
        %~ \highlightbf{Strategy} &  \\
        %~ \hline
        %~ \highlightbf{Command} & same considerations 
                                %~ for~\hyperlink{tab:command-relation-with-others}{Command} \\
        %~ \hline      
        %~ \highlightbf{Decorator} & changes the "skin" whereas \highlightit{Strategy} changes the "guts"\\
        %~ \hline      
        %~ \highlightbf{Template Method} & is based on inheritance and is
                                        %~ useful to change parts of an 
                                        %~ algorithm in subclasses, 
                                        %~ whereas \highlightit{Strategy} 
                                        %~ is based on composition and 
                                        %~ alters the object's behavior 
                                        %~ with different  \\
        %~ \hline      
        %~ \highlightbf{State} & is based on composition as well \\
        %~ \hline      
    %~ \end{tabularx}
%~ \end{table}

\newpage
%%% %%% TEMPLATE METHOD %%% %%% %%%
\subsection{Template Method}
\begin{table}[h!t]
    \begin{tabularx}{\textwidth}{|l|X|}
        \hline
        \multirow{1}{*}{\textbf{What for}} & \highlightit{defining the skeleton of 
                                             an algorithm in a class} \\
                                           & and letting subclasses 
                                             override specific steps\\
        \hline
        \multirow{1}{*}{\textbf{How to}} & based on \highlightit{inheritance}\\
        \hline
        \textbf{When needed} & $\bullet$ client needs to extend only 
                                         specific steps of an algorithm \\
                             & $\bullet$ managing better serveral 
                                         classes containing almost 
                                         identical algorithms\\
        \hline
    \end{tabularx}
\end{table}

\toclessnparagraph{Structure}
\vspace{-1.5cm}
\begin{table}[h!t]
    \begin{center}
    \begingroup
    \begin{figure}[H]
        \begin{center}
    
    \begin{tikzpicture}%[show background grid]
        
        \begin{abstractclass}[text width=3.5cm]{ABCClass}{0,0}
            \attributett{...}
            \operationtt{+ templateMethod()}
            \operationtt{+ primitive1()}
            \operationtt{+ primitive2()}
            \operationtt{+ primitive3()}
            \operationtt{+ primitive4()}
        \end{abstractclass}
        \begin{class}[text width=3.5cm]{ClassA}{-2,-6}
            \implement{ABCClass}
            \attributett{...}
            \operationtt{+ prititive3()}
            \operationtt{+ primitive4()}            
        \end{class}
        \begin{class}[text width=3.5cm]{ClassB}{2,-6}
            \implement{ABCClass}
            \attributett{...}
            \operationtt{+ primitive1()}
            \operationtt{+ primitive2()}
            \operationtt{+ primitive3()}
            \operationtt{+ primitive4()}            
        \end{class}
        
    \end{tikzpicture}
    
\end{center}
        
    \end{figure}
%~ \begin{itemize}
    %~ \item \textbf{Abstract Class} declares the steps of an algorithm and 
    %~ the \textbf{templateMethod}  used to call them in a specific order. The steps may both be 
    %~ declared as absrtact or implement some actual routines.
    %~ \item \textbf{Concrete Classes} override all of the steps, but not 
    %~ the \textbf{templateMethod} itself.
%~ \end{itemize}
    \endgroup
    \begin{tabularx}{\textwidth}{|l|X|}
        \hline
        \highlightbf{ABCClass} &  declares the steps of an algorithm and 
                                  the \highlighttt{templateMethod} used to call 
                                  them in a specific order. The steps 
                                  may both be declared as absrtact or 
                                  implement some actual routines.\\
        \hline
        \highlightbf{ClassA}, \highlightbf{ClassB} & override all of the steps, 
                                                     but not the 
                                                     \highlighttt{templateMethod} itself.\\
        \hline
    \end{tabularx}    
    \caption{Strategy Class Diagram.}
    \label{tab:strategy-class-diagram}     
    \end{center}
\end{table}  

\newpage
\toclessnparagraph{Pros and Cons}
\begin{table}[h!t]
    \begin{tabularx}{\textwidth}{|l|X|}
        \hline
        \textbf{Pros} & $\bullet$ overriding part of algorithms without 
                                  affecting the rest \\
                      & $\bullet$ less code duplication\\
        \hline
        \textbf{Cons} & $\bullet$ limitations due to the skeleton 
                        of the algorithm\\
                      & $\bullet$ \highlightit{Liskov Substitution Principle}\\
                      & $\bullet$ maintainance can get harder as the steps increase\\
        \hline      
    \end{tabularx}
\end{table}

%~ \toclessnparagraph{Relation with other Design Patterns}
%~ \begin{table}[h!t]
    %~ \begin{tabularx}{\textwidth}{|l|X|}
        %~ \hline
        %~ \highlightbf{Factory Method} & a specialization of \highlightit{Template Method} and may be a step in a large one of the latter\\
        %~ \hline
        %~ \highlightbf{Strategy} & differently from \highlightit{Template Method}, is based on \highlightit{composition} and works on object level, at run time \\
        %~ \hline      
    %~ \end{tabularx}
%~ \end{table}

\newpage
%%% %%% VISITOR %%% %%% %%%
\subsection{Visitor}
\begin{table}[h!t]
    \begin{tabularx}{\textwidth}{|l|X|}
        \hline
        \multirow{1}{*}{\textbf{What for}} & \highlightit{separating algorithms 
                                             from the objects} they 
                                             operate on\\
        \hline
        \multirow{1}{*}{\textbf{How to}} & \\
        \hline
        \textbf{When needed} & $\bullet$ performing an operation on all elements 
                               of a complex data structure\\
                             & $\bullet$ cleaning up the business 
                               logic extracting other behaviors\\
                             & $\bullet$ separating behaviors that 
                               make sense only for specific classes\\
        \hline
    \end{tabularx}
\end{table}

\toclessnparagraph{Structure}
\vspace{-1.5cm}
\begin{table}[h!t]
    \begin{center}
    \begingroup
    \begin{figure}[H]
        \begin{center}
\scalebox{0.9}
{    
    \begin{tikzpicture}%[show background grid]
        
        \begin{interface}[text width=3.8cm]{Visitor}{0,0}
            \operationtt{+ visit(e: ElementA)}
            \operationtt{+ visit(e: ElementB)}
        \end{interface}
        \begin{class}[text width=3.8cm]{VisitorA}{-2.5,-4}
            \implement{Visitor}
            \operationtt{+ visit(e: ElementA)}
            \operationtt{+ visit(e: ElementB)}
        \end{class}
        \begin{class}[text width=3.8cm]{VisitorB}{2.5,-4}
            \implement{Visitor}
            \operationtt{+ visit(e: ElementA)}
            \operationtt{+ visit(e: ElementB)}        
        \end{class}
        
        \begin{interface}[text width=4.5cm]{Element}{10,0}
            \operationtt{+ accept(v: Visitor)}
        \end{interface}
        \begin{class}[text width=4.5cm]{ElementA}{7.5,-3}
            \implement{Element}
            \attributett{...}
            \operationtt{+ featureA()}
            \operationtt{+ accept(v: Visitor)}
        \end{class}
        \begin{class}[text width=4.5cm]{ElementB}{12.5,-3}
            \implement{Element}
            \attributett{...}
            \operationtt{+ featureB()}
            \operationtt{+ accept(v: Visitor)}
        \end{class}

        \begin{class}[text width=3cm]{Client}{5,-7}
        \end{class}
        
    \end{tikzpicture}
}    
\end{center}

      
    \end{figure}
    \endgroup

    \begin{tabularx}{\textwidth}{|l|X|}
        \hline
        \highlightbf{Visitor} & interface declares a set of visiting methods, 
    each with a different Concrete Element taken as argument.\\
        \hline
        \highlightbf{VisitorA}, \highlightbf{VisitorB} & implement the different behaviors for 
                                                         each Concrete Element. \\
        \hline
        \highlightbf{Element} & interface declares a method to accept the 
                                Visitor interface as argument. \\
        \hline
        \highlightbf{ElementA}, \highlightbf{ElementB} & implements the acceptance method, 
    whose purpose is the redirection the call to the proper visitor's method 
    for the current element class. This method must be overridden in any 
    subclass even though it was already developed in the base one.\\
        \hline
        \highlightbf{Client} & usually needs to work on a collection or a similar complex 
    object whose type is ignored by the Client, since it works with 
    the abstract interface.\\
        \hline
    \end{tabularx}    
    \caption{Strategy Class Diagram.}
    \label{tab:strategy-class-diagram}     
    \end{center}    
\end{table} 
   
\toclessnparagraph{Pros and Cons}
\begin{table}[h!t]
    \begin{tabularx}{\textwidth}{|l|X|}
        \hline
        \textbf{Pros} & $\bullet$ \highlightit{Open/Close Principle}\\
                      & $\bullet$ \highlightit{Single Responsibility Principle}\\
                      & $\bullet$ can store information while working 
                                  with different objects\\
        \hline
        \textbf{Cons} & $\bullet$ adding/removing a class requires 
                        all visitors to get updated \\
                      & $\bullet$ access grants to private data and 
                                  methods of objects they are 
                                  working on \\
        \hline      
    \end{tabularx}
\end{table}

%~ \toclessnparagraph{Relation with other Design Patterns}
%~ \begin{table}[h!t]
    %~ \begin{tabularx}{\textwidth}{|l|X|}
        %~ \hline
        %~ \highlightbf{Command} & \highlight{Visitor} can be treated as a powerful version of \highlightit{Command}
                                %~ letting the user execute operation over various objects of different classes\\
        %~ \hline
        %~ \highlightbf{Composite} & \highlight{Visitor} can execute an operation over an entire \highlight{Composite} tree \\
        %~ \hline      
        %~ \highlightbf{Iterator} &  along with \highlightit{Visitor} let users traverse a complex data structure and execute 
                                  %~ some operation over its elements, even if belonging to different classes\\
        %~ \hline      
    %~ \end{tabularx}
%~ \end{table}

\newpage 
%%% STRUCTURAL PATTERNS %%%
\section{Structural Patterns}
%%% %%% ADAPTER %%% %%% %%%
\subsection{Adapter}
\begin{table}[h!t]
    \begin{tabularx}{\textwidth}{|l|X|}
        \hline
        \multirow{1}{*}{\textbf{What for}} & allowing \highlightit{objects with 
        incompatible interfaces to collaborate}\\
        \hline
        \multirow{1}{*}{\textbf{How to}} & \\
        \hline
        \textbf{When needed} & $\bullet$ a class' interface is 
                               incompatible with the rest of the 
                               code\\
                             & $\bullet$ reusing sublcasses that lack 
                                         some common functionalities 
                                         that can't be added in the 
                                         superclass\\
        \hline
    \end{tabularx}
\end{table}
%~ \textbf{Adapter} allows the collaboration among objects with incompatible interfaces.
%~ This is achieved through the implementation of an Adapter which make 
%~ one interface understandable by an object of a different class.
%~ The adatper gets an interface compatible with one existing object.
%~ Through this interface, the objects can call the adapter's methods.
%~ Upon receiving a call, the adapter passes the request to the second object,
%~ in a format and order it could understand.
%~ An Adapter could work in both directions.

%~ Since the Client interacts with the adapter through the client interface, 
%~ it doesn't get coupled. Therefore, different adapters can be added to the program 
%~ without breaking any existing code.

\toclessnparagraph{Structure}
\vspace{-1.5cm}
\begin{table}[h!t]
    \begin{center}
    \begingroup
    \begin{figure}[H]
    \begin{center}
    
    \begin{tikzpicture}[show background grid]
        
        \begin{abstractclass}[text width=3cm]{Client Interface}{0,0}
            \operationtt{+ method(data)}
        \end{abstractclass}
        \begin{class}[text width=3cm]{Adapter}{0,-3}
            \attributett{- adaptee: Service}
            \operationtt{+ mwthod(data)}
        \end{class}
        
        \begin{class}[text width=5cm]{Service}{6,-3}
            \attributett{...}
            \operationtt{serviceMethod(specialData)}
        \end{class}
        
        \begin{class}[text width=2cm]{Client}{-4,-0.3}
            
        \end{class}
        
    \end{tikzpicture}
    
\end{center}
      
    \end{figure}
    \endgroup
    \begin{tabularx}{\textwidth}{|l|X|}
        \hline
        \highlightbf{Client} & the business logic of the program\\
        \hline
        \highlightbf{ClientInterface} & the protocol for a class to 
    be able to collaborate with it.\\
        \hline
        \highlightbf{Service} & usually a 3rd party or legacy software 
    with an interface incopatible with the Client.\\
        \hline
        \highlightbf{Adapter} & a class able to work with both the client and 
    the service. It receives calls from the client and translates 
    into calls to a wrapped service object.\\
        \hline
    \end{tabularx}    
    \caption{Strategy Class Diagram.}
    \label{tab:strategy-class-diagram}     
    \end{center}      
\end{table}

\toclessnparagraph{Pros and Cons}
\begin{table}[h!t]
    \begin{tabularx}{\textwidth}{|l|X|}
        \hline
        \textbf{Pros} & $\bullet$ \highlightit{Single Responsibility Principle}\\
                      & $\bullet$ \highlightit{Open/Close Principle}\\
        \hline
        \textbf{Cons} & $\bullet$ sometimes it's simpler to change 
                                  the service class in order to match 
                                  than adding a new set of interfaces\\
        \hline      
    \end{tabularx}
\end{table}

%~ \newpage
%~ \toclessnparagraph{Relation with other Design Patterns}    
%~ \begin{table}[h!t]
    %~ \begin{tabularx}{\textwidth}{|l|X|}
        %~ \hline
        %~ \highlightbf{Bridge} & usually designed in advance to develop 
                               %~ independent parts of an application, 
                               %~ whereas \highlightit{Adapter} makes
                               %~ incompatible classes work together \\
        %~ \hline
        %~ \highlightbf{Decorator} & supports \highlightit{recursive} 
                                  %~ composition and can let the interface 
                                  %~ stay the same or be extended \\
        %~ \hline      
        %~ \highlightbf{Proxy} & differently from \highlightit{Adapter} 
                              %~ makes users access another object using 
                              %~ the same interface\\
        %~ \hline      
        %~ \highlightbf{Facade} & defines a new interface for an already 
                               %~ existing object and can work with an 
                               %~ entire subsystem of objects\\
        %~ \hline      
        %~ \highlightbf{Bridge} & have similar structures, 
                               %~ based on \highlightit{composition}, 
                               %~ delegating work to  \\
        %~ \highlightbf{State} &  other objects.  \\
        %~ \highlightbf{Strategy} & They solve different problems though. \\
        %~ \hline      
    %~ \end{tabularx}
%~ \end{table}

\newpage
%%% %%% BRIDGE %%% %%% %%%
\subsection{Bridge}
\begin{table}[h!t]
    \begin{tabularx}{\textwidth}{|l|X|}
        \hline
        \multirow{1}{*}{\textbf{What for}} & \highlightit{splitting a large class 
                                             or a set of related classes 
                                             into two hierarchies}, an 
                                             abstracion and an implementation, 
                                             developed independently from each other\\
        \hline
        \multirow{1}{*}{\textbf{How to}} & \\
        \hline
        \textbf{When needed} & $\bullet$ dividing and organizing a 
                                         monolithic class that has 
                                         several variants of some 
                                         functionality \\
                             & $\bullet$ extending a class in several 
                                         independent dimensions\\
                             & $\bullet$ switching implementation at 
                                         runtime is a requirement\\
        \hline
    \end{tabularx}
\end{table}
%~ \textbf{Bridge} splits a large class or a set of closely related classes 
%~ into two separate hierarchies which can be developed independently of each other.
%~ This is achieved by switching from inheritance to object composition.
%~ One of the dimensions is extracted into a separate class hierarchy, 
%~ therefore the original classes will reference an object of the new hierarchy
%~ instead of having all its states and behaviors coded within the class.
%~ The original class implement a reference field for each "extracted-dimension" classes 
%~ This extraction makes a class able to delegate any 'dimension-related' work
%~ to the linked 'dimension-related' object.

%~ \begin{itemize}
    %~ \item \textbf{Abstraction} provides high-level interface, and relies 
    %~ on the Implementation objects to perform the actual work.
    %~ \item \textbf{Implementation} is the interface shared by all the 
    %~ concrete implementations. Usually the abstraction declares some 
    %~ complex behavior that relies on primitive operations declared 
    %~ by the implementation.
    %~ \item \textbf{Concrete Implementations1} implement specific code.
    %~ \item \textbf{Refined Abstraction} are optional and might help 
    %~ providing variants of control logic.
    %~ \item \textbf{Client} interacts with the abstraction, after linking 
    %~ it with one of the implementation objects.
%~ \end{itemize}

\toclessnparagraph{Structure}
\vspace{-1.5cm}
\begin{table}[h!t]
    \begin{center}
    \begingroup
    \begin{figure}[H]
            \begin{tikzpicture}[show background grid]
        \begin{class}[text width=5cm]{Abstraction}{0,0}
            \attributett{- impl: Implementation}
            \operationtt{+ feature1()}
            \operationtt{+ feature2()}
        \end{class}
        
        \begin{class}[text width=3.5cm]{RefinedAbstraction}{0,-4}
            \inherit{Abstraction}
            \attributett{...}
            \operationtt{+ featureN()}
        \end{class}
        
        \begin{abstractclass}[text width=4cm]{Implementation}{8,0.5}
            \operationtt{+ method1()}
            \operationtt{+ method2()}
            \operationtt{+ method3()}
        \end{abstractclass}
        
        \begin{class}[text width=3cm]{ConcreteImplA}{6,-4}
        \implement{Implementation}\end{class}

        \begin{class}[text width=3cm]{ConcreteImplB}{10,-4}
        \implement{Implementation}\end{class}
        
        \aggregation{Abstraction}{}{}{Implementation}
    \end{tikzpicture}

    \end{figure}
    \endgroup
    \begin{tabularx}{\textwidth}{|l|X|}
        \hline
        \highlightbf{Abstraction} & high-level interface, and relies 
    on the Implementation objects to perform the actual work.\\
        \hline
        \highlightbf{RefinedAbstraction} & optional and might help 
    providing variants of control logic.\\
        \hline
        \highlightbf{Implementation} & interface shared by all the 
    concrete implementations. Usually the abstraction declares some 
    complex behavior that relies on primitive operations declared 
    by the implementation.\\
        \hline
        \highlightbf{ImplementationA}, &  implement specific code\\
        \highlightbf{ImplementationB} & \\
        \hline
        \highlightbf{Client} & interacts with the abstraction, after linking 
    it with one of the implementation objects.\\
        \hline
    \end{tabularx}    
    \caption{Bridge Class Diagram.}
    \label{tab:bridge-class-diagram}     
    \end{center}     
\end{table}

\newpage
\toclessnparagraph{Pros and Cons}
\begin{table}[h!t]
    \begin{tabularx}{\textwidth}{|l|X|}
        \hline
        \textbf{Pros} & $\bullet$ platform-independent classes are easier to 
                         create\\
                      & $\bullet$ client isn't exposed to platform details 
                         and work with high-level abstractions\\
                      & $\bullet$ \highlightbf{Open/Close Principle} \\
                      & $\bullet$ \highlightbf{Single Responsibility Principle} \\
        \hline
        \textbf{Cons} & an highly cohesive class might require code 
                        getting more complicated   \\
        \hline      
    \end{tabularx}
\end{table}

%~ \toclessnparagraph{Relation with other Design Patterns}    
%~ \begin{table}[h!t]
    %~ \begin{tabularx}{\textwidth}{|l|X|}
        %~ \hline
        %~ \highlightbf{Adapter} & is usually implemented for and 
                                %~ already existing app, in order to 
                                %~ make incompatible classes work together \\
        %~ \hline
        %~ \highlightbf{Bridge} & have similar structures, 
                               %~ based on \highlightit{composition}, 
                               %~ delegating work \\
        %~ \highlightbf{State} &  to other objects.  \\
        %~ \highlightbf{Strategy} & They solve different problems though. \\            \hline
        %~ \highlightbf{Abstract Factory} & can be used along with 
                                         %~ \highlightit{Bridge}, 
                                         %~ encapsulating the abstraction
                                         %~ the latter may be bounded 
                                         %~ to work with.  \\
        %~ \hline
        %~ \highlightbf{Builder} & can be combined with \highlightit{Builder}, 
                                %~ with the \highlightit{director} class 
                                %~ playing the role of the abstraction 
                                %~ and different builder acting as 
                                %~ implementations \\
        %~ \hline      
    %~ \end{tabularx}
%~ \end{table}
    
\newpage
%%% %%% COMPOSITE %%% %%% %%%
\subsection{Composite}
\begin{table}[h!t]
    \begin{tabularx}{\textwidth}{|l|X|}
        \hline
        \multirow{1}{*}{\textbf{What for}} & \highlightit{composing objects 
                                             into tree structures} in order 
                                             to work with these structures 
                                             as if they were individual 
                                             objects\\
        \hline
        \multirow{1}{*}{\textbf{How to}} & both leaf and composite 
                                           objects implement the same 
                                           component interface, the 
                                           former working as a single 
                                           entity and the latter 
                                           possibly containing 
                                           sub-elements, routines to 
                                           add and remove them, and 
                                           a routine to delegate work
                                           to them \\
        \hline
        \textbf{When needed} & $\bullet$ implementing a tree-like 
                                         data structures \\
                             & $\bullet$ clients need to interact 
                                         uniformly with simple and 
                                         complex objects \\
        \hline
    \end{tabularx}
\end{table}
%~ \textbf{Composite} after enabling composition of objects into tree structures, 
%~ this pattern let clients work with these structures as they were individual objects.
%~ This is achieved extending the class



\toclessnparagraph{Structure}
\vspace{-1.5cm}
\begin{table}[h!t]
    \begin{center}
    \begingroup
    \begin{figure}[H]
        \centering    
    \begin{tikzpicture}[show background grid]
    
        \begin{abstractclass}[text width=3cm]{ABCComponent}{0,0}
            \operationtt{+ execute()}
        \end{abstractclass}
        \begin{class}[text width=3cm]{Leaf}{0,-3.5}
            \attributett{...}
            \operationtt{+ execute()}
        \end{class}
        \begin{class}[text width=5.5cm]{Composite}{6,-3.5}
            \attributett{- children: Component[]}
            \operationtt{+ add(c: Component)}
            \operationtt{+ remove(c: Component)}
            \operationtt{+ getChildren(): Component[]}
            \operationtt{+ execute()}
        \end{class}

        \begin{class}[text width=3cm]{Client}{0,2}
        \end{class}
            
    \end{tikzpicture}
       
    \end{figure}
    \endgroup
    \begin{tabularx}{\textwidth}{|l|X|}
        \hline
        \highlightbf{Component} & interface contains the common operations 
    for both simple and complex elements of the tree.\\
        \hline
        \highlightbf{Leaf} & a basic element of a tree and doesn't contain 
    any other sub-element. Usually is the part of the system which ends up 
    doing the work.\\
        \hline
        \highlightbf{Composite} & an element that has sub-elements, 
    which means leafs and containers as well. It interacts with sub-elements 
    through their interfaces.\\
        \hline
        \highlightbf{Client} & can work with both simple and complex elements 
    through the component interface.\\
        \hline    
    \end{tabularx}
    \caption{Composite Class Diagram.}
    \label{tab:composite-class-diagram} 
    \end{center}     
\end{table}

\newpage
\toclessnparagraph{Pros and Cons}
\begin{table}[h!t]
    \begin{tabularx}{\textwidth}{|l|X|}
        \hline
        \textbf{Pros} & $\bullet$ polymorphism and recursion help 
                        work easier with complex structures \\
                      & $\bullet$ \highlightbf{Open/Close Principle} \\
        \hline
        \textbf{Cons} & a common interface can be difficult when 
                        functionalities of single elements differe 
                        too much from the ones of composite ones \\
        \hline      
    \end{tabularx}
\end{table}
%~ \newpage
%~ \toclessnparagraph{Relation with other Design Patterns}    
%~ \begin{table}[h!t]
    %~ \begin{tabularx}{\textwidth}{|l|X|}
        %~ \hline
        %~ \highlightbf{Builder} & can be used to construct a complex composite 
                           %~ object step to step \\
        %~ \hline
        %~ \highlightbf{Chain of Responsibility} & can be used along with 
        %~ composite leaf nodes to build a tree of responsibilities \\
        %~ \hline      
        %~ \highlightbf{Iterator}s & can be used to traverse composite trees \\
        %~ \hline      
        %~ \highlightbf{Visitor} & can be used to execute an operation over an 
                           %~ entire composite tree \\
        %~ \hline      
        %~ \highlightbf{Flyweight}s & can be used to implement shared leaf nodes \\
        %~ \hline      
        %~ \highlightbf{Decorator} & is similar to \textit{Composite} with 
                             %~ only one child component, but it adds 
                             %~ functionalities to the wrapped object\\
        %~ \hline      
        %~ \highlightbf{Prototype} & can help designs with heavy use of 
                             %~ \textit{Composite} and \textit{Decorator}, 
                             %~ making possible to clone complex structures 
                             %~ instead of rebuilding from scratch  \\
        %~ \hline      
    %~ \end{tabularx}
%~ \end{table}

    
\newpage
%%% %%% DECORATOR %%% %%% %%%
\subsection{Decorator}
\begin{table}[h!t]
    \begin{tabularx}{\textwidth}{|l|X|}
        \hline
        \multirow{1}{*}{\textbf{What for}} & \highlightit{attach new behaviors to 
                                             objects}, placing them in special 
                                             wrappers\\
        \hline
        \multirow{1}{*}{\textbf{How to}} & objects are placed inside 
                                           special wrapper objects 
                                           that contain the behaviors\\
        \hline
        \textbf{When needed} & $\bullet$ extra behavior needs to be assigned to 
                               objects at runtime\\
                             & $\bullet$ extending object's behavior when 
                               inheritance can not be used \\
                             & \quad (for example when a class is 
                               declared as \texttt{final}\\
        \hline
    \end{tabularx}
\end{table}
%~ \textbf{Decorator} attaches new behaviors to objects by placing these objects inside
%~ special wrapper objects containing the desired behavior.

%~ Multiple layers of decorations are possible as long as the Client 
%~ interacts with the object through the component.
    
\toclessnparagraph{Structure}
\vspace{-1.5cm}
\begin{table}[h!t]
    \begin{center}
    \begingroup
    \begin{figure}[H]
        \begin{center}
    
    \begin{tikzpicture}[show background grid]
        
        \begin{abstractclass}[text width=3cm]{ABCComponent}{0,0}
            \operationtt{+ execute()}
        \end{abstractclass}
        
        \begin{class}[text width=3cm]{ComponentA}{0,-3}
            \implement{ABCComponent}
            \attributett{...}
            \operationtt{+ execute()}
        \end{class}
        
        \begin{abstractclass}[text width=5.5cm]{ABCDecorator}{6,-3}
            \attributett{- wrappee: Component}
            \operationtt{+ ABCDecorator(c: Component)}
            \operationtt{+ execute()}
        \end{abstractclass}
        \begin{class}[text width=3cm]{DecoratorA}{4,-7}
            \implement{ABCDecorator}
            \attribute{...}
            \operation{+ execute()}
            \operation{+ extra()}
        \end{class}
        \begin{class}[text width=3cm]{DecoratorB}{8,-7}
            \implement{ABCDecorator}
            \attribute{...}
            \operation{+ execute()}
            \operation{+ extra()}
        \end{class}
        
        \begin{class}[text width=3cm]{Client}{0,2}
        \end{class}
        
    \end{tikzpicture}
    
\end{center}
      
    \end{figure}
    \endgroup
    \begin{tabularx}{\textwidth}{|l|X|}
        \hline
        \highlightbf{Component} & interface common for both wrapper and 
    wrapped objects.\\
        \hline
        \highlightbf{ComponentA} &  class being wrapped and whose 
    basic behavior can be altered by decorators.\\
        \hline
        \highlightbf{Decorator} & has a reference of the wrapped object, 
    whose referenced type matches the Component interface. The Base Decorator 
    delegates all operations to the wrapped object.\\
        \hline
        \highlightbf{DecoratorA}, & define behaviors which can be added 
    to Components dynamically. \\
        \highlightbf{DecoratorB} & The concrete decorator overrides the 
                                   base behavior and execute the variant either before of after calling the 
                                   parent method.\\
        \hline
    \end{tabularx}    
    \caption{Decorator Class Diagram.}
    \label{tab:decorator-class-diagram} 
    \end{center}
\end{table}

\newpage
\toclessnparagraph{Pros and Cons}
\begin{table}[h!t]
    \begin{tabularx}{\textwidth}{|l|X|}
        \hline
        \textbf{Pros} & no other classes are needed to extend an 
                        object's behavior\\
                      & adding and removing responsibilities can be 
                        achieved at runtime\\
                      & ?? \\
        \hline
        \textbf{Cons} & $\bullet$ removing wrappers from the stack 
                        can be difficult\\
                      & $\bullet$ implementing a decorator in a way 
                        its behavior does not depend on the order 
                        in the stack can be hard\\
                      & $\bullet$ code can get ugly when the layers 
                        increase in number\\
        \hline      
    \end{tabularx}
\end{table}

%~ \toclessnparagraph{Relation with other Design Patterns}    
 %~ \begin{table}[h!t]
    %~ \begin{tabularx}{\textwidth}{|l|X|}
        %~ \hline
        %~ \highlightbf{Adapter} &  \\
        %~ \hline
        %~ \highlightbf{Chain of Responsibility} & has a similar class 
                                                %~ structure that relies 
                                                %~ on recursive composition, 
                                                %~ but its handlers can 
                                                %~ operate independently and 
                                                %~ stop passing the request 
                                                %~ at any point, whereas 
                                                %~ \highlightit{Decorator} 
                                                %~ is not allowed to break 
                                                %~ the flow 
                                                 %~ \\
        %~ \hline
        %~ \highlightbf{Composite} & relies on recursive composition \\
        %~ \hline
        %~ \highlightbf{Prototype} &  \\
        %~ \hline
        %~ \highlightbf{Startegy} & change the internal behaviour whereas 
                                 %~ \highlightit{Decorator} changes the exterior \\
        %~ \hline
        %~ \highlightbf{Proxy} & similar structure based on composition, 
                              %~ but different intent: \highlightit{Proxy} 
                              %~ usually manages the life cycle of its 
                              %~ service object, whereas the client 
                              %~ controls the composition of 
                              %~ \highlightit{Decorator} \\
        %~ \hline
    %~ \end{tabularx}
%~ \end{table}
       
\newpage
%%% %%% FACADE %%% %%% %%%
\subsection{Facade}
\begin{table}[h!t]
    \begin{tabularx}{\textwidth}{|l|X|}
        \hline
        \multirow{1}{*}{\textbf{What for}} & providing users with a 
                                             \highlightit{simplified 
                                             interface to any complex set of} \highlightit{classes}\\
        \hline
        \multirow{1}{*}{\textbf{How to}} & a facade class provides a 
                                           simple interface to a 
                                           complex subsystem\\
        \hline
        \textbf{When needed} & $\bullet$ limited but straightforward 
                               interface to a complex subsystem is \\ 
                             & \quad needed\\
                             & $\bullet$ subsystem needs to be 
                               structured into layers\\
        \hline
    \end{tabularx}
\end{table}
%~ \textbf{Facade} consists of a simplified interface to a complex subsystem 
%~ which contains lots of moving parts. Functionalities provided by the Facade 
%~ could be reduced to the subset the client is interested in.
%~ \begin{itemize}
    %~ \item \textbf{Facade} implements a convenient access to a particular 
    %~ part of the subsystem's functionality and knows where to direct 
    %~ the client's request and how to operate all the system's parts.
    %~ \item \textbf{Additional Facade} might help reducing the tasks 
    %~ managed by a single Facade. It can be used by the Client and 
    %~ by the other facades as well.
    %~ \item \textbf{Complex Subsystem} Subsytem class are various object 
    %~ which need to be studied deeply in order to manage them properly. 
    %~ These classes work together and are unaware of the Facade's existence.
    %~ \item \textbf{Client} calls the subsystem classes through the Facade only. 
%~ \end{itemize}
    
\toclessnparagraph{Structure}
\vspace{-1.5cm}
\begin{table}[h!t]
    \begin{center}
    \begingroup
    \begin{figure}[H]
        \begin{center}
    
    
    \tikzfig{facade_design_pattern}
\begin{tikzpicture}%[show background grid]
        
        \begin{class}[text width=5cm]{Facade}{0,0}
            \attributett{- linksToSubSysObjs}
            \attributett{- optionalAdditionalFacde}
            \operationtt{+ subSysOperation()}
        \end{class}
        
        \begin{class}[text width=3cm]{SubSystemA}{-3.5,-3.5}
        \end{class}
        \begin{class}[text width=3cm]{SubSystemB}{0,-3.5}
        \end{class}
        \begin{class}[text width=3cm]{SubSystemC}{3.5,-3.5}
        \end{class}
        
        
        \begin{class}[text width=4.5cm]{AdditionalFacade}{8.5,0}
            \operationtt{+ anotherSubSysOperation}
        \end{class}
        
        %~ \begin{class}[text width=4.5cm]{AdditionalSubSys}{8.5,-3.5}
        %~ \end{class}
        
        
        \begin{class}[text width=3cm]{Client}{0,2}
        \end{class}
        
        \draw [ - >] ( Client)  -- ( Facade) ;                                   
        \draw [ - >] ( Facade)  -- ( AdditionalFacade) ;                                   
        
        \draw[dashed, - >] (Facade) -- (SubSystemA);
        \draw[dashed, - >] (Facade) -- (SubSystemB);
        \draw[dashed, - >] (Facade) -- (SubSystemC);
        
    \end{tikzpicture}
    

\end{center}
       
    \end{figure}
    \endgroup
    \begin{tabularx}{\textwidth}{|l|X|}
        \hline
        \highlightbf{Facade} & implements a convenient access to a particular 
    part of the subsystem's functionality and knows where to direct 
    the client's request and how to operate all the system's parts.\\
        \hline
        \highlightbf{AdditionalFacade} & might help reducing the tasks 
    managed by a single Facade. It can be used by the Client and 
    by the other facades as well.\\
        \hline
        \highlightbf{SubSystemA}, \dots, & Subsytem class are various object 
    which need to be studied \\
        \highlightbf{SubSystemC} & deeply in order to manage them properly. 
    These classes work together and are unaware of the Facade's existence.\\
        \hline
        \highlightbf{Client} & calls the subsystem classes through the Facade only\\
        \hline
    \end{tabularx}    
    \caption{Facade Class Diagram.}
    \label{tab:facade-class-diagram} 
    \end{center}
\end{table}

\toclessnparagraph{Pros and Cons}
\begin{table}[h!t]
    \begin{tabularx}{\textwidth}{|l|X|}
        \hline
        \textbf{Pros} & client code is isolated from the complexity of a subsystem\\
        \hline
        \textbf{Cons} & desire to couple a class to a facade can 
                        overwhelm the code base \\
        \hline      
    \end{tabularx}
\end{table}

%~ \newpage
%~ \toclessnparagraph{Relation with other Design Patterns}    
 %~ \begin{table}[h!t]
    %~ \begin{tabularx}{\textwidth}{|l|X|}
        %~ \hline
        %~ \highlightbf{Adapter} & instead of defining a new interface as 
                           %~ \textit{Facade} for several subsystems, 
                           %~ it tries to adapt to an existing one \\
        %~ \hline      
        %~ \highlightbf{Abstract Factory} & alternative in case it is only 
                                    %~ needed to hide the objects 
                                    %~ creation details from the client 
                                    %~ code \\
        %~ \hline      
        %~ \highlightbf{Flyweight} & instead of representing an entire 
                             %~ subsystem, it makes it easier to make 
                             %~ lots of little objects \\
        %~ \hline      
        %~ \highlightbf{Mediator} & it also tries to organize collaboration 
                            %~ between lots of tightly coupled class \\
        %~ \hline      
        %~ \highlightbf{Singleton} & they are equivalent in case a single facade 
                             %~ object is needed \\
        %~ \hline      
        %~ \highlightbf{Proxy} & it also buffers a complex entity and 
                         %~ initializes it, but the interface is not the 
                         %~ same as its service object\\
        %~ \hline      
    %~ \end{tabularx}
%~ \end{table}
        
\newpage
%%% %%% FLYWEIGHT %%% %%% %%%
\subsection{Flyweight}
\begin{table}[h!t]
    \begin{tabularx}{\textwidth}{|l|X|}
        \hline
        \multirow{1}{*}{\textbf{What for}} & reducing RAM consumption 
                                             by \highlightit{sharing common parts 
                                             of state} \highlightit{between multiple 
                                             objects}, instead of creating 
                                             copies of same data in 
                                             each object\\
        \hline
        \multirow{1}{*}{\textbf{How to}} & after extracting the 
                                           \highlightit{intrinsic state} 
                                           (invariant, context-independent 
                                           and shareable) an interface 
                                           to the 
                                           \highlightit{extrinsic state} 
                                           (varinat, \\
                                           & context-dependent and 
                                           unshareable) is implemented\\
        \hline
        \textbf{When needed} & $\bullet$ dealing with large numbers 
                               of objects with simple repeated \\ 
                             & \quad elements that would use a large 
                               amount of memory if \\
                             & \quad individually stored\\
                             %~ & $\bullet$  \\
        \hline
    \end{tabularx}
\end{table}

\toclessnparagraph{Structure}
\vspace{-1.5cm}
\begin{table}[h!t]
    \begin{center}
    \begingroup
    \begin{figure}[H]
        \begin{center}
    
    \begin{tikzpicture}%[show background grid]

    \begin{class}[text width=5.5cm]{FlyweightFactry}{0,0}
        \attributett{- cache: Flyweight[]}
        \operationtt{+ getFlyweight(extrinsicState)}
    \end{class}

    \begin{interface}[text width=5.5cm]{Flyweight}{0,-3.5}
        \attributett{- cache: Flyweight[]}
        \operationtt{+ getFlyWeight(extrinsicState)}
        \operationtt{+ operation(extrinsicState)}
    \end{interface}
    \begin{class}[text width=5cm]{SharedFlyweight}{-3,-8.5}
        \implement{Flyweight}
        \attributett{- intinsicState}
        \operationtt{+ operation(extrinsicState)}
    \end{class}
    \begin{class}[text width=5cm]{UnsharedFlyweight}{3,-8.5}
        \implement{Flyweight}  
        \operationtt{+ operation(extrinsicState)}      
    \end{class}

    \begin{class}[text width=2cm]{Client}{-6,-0.3}
    \end{class}



    \end{tikzpicture}
    
\end{center}
      
    \end{figure}
    \endgroup
%~ \begin{itemize}
    %~ \item \textbf{Flyweight} contains the portion of the original object's state 
    %~ that can be shared between multiple objects. The state stored 
    %~ inside a flyweight is called \textbf{\textit{intrinsic}}; the state 
    %~ passed to the flyweight's methods is called \textbf{\textit{extrinsic}}. 
    %~ The same flyweight object can be used in many contexts.
    %~ \item \textbf{Context} contains the extrinsic state, unique across all original 
    %~ objects. It represents the full state when paired with one of the 
    %~ flyweight objects.
    %~ \item \textbf{Flyweight Factory} manages a pool of existing flyweights 
    %~ freeing the Client from the task of creating. The Flyweight Factory 
    %~ receives a call to retrieve a specific flyweight, looks over previously 
    %~ created flyweight and either returns one of these or creates 
    %~ a ne wone in case nothing matches the search criteria.
    %~ \item \textbf{Client} sees the flyweight as a template object 
    %~ which can be configured at runtime by passing some contextual 
    %~ data into parameters of its methods.
%~ \end{itemize}
    \begin{tabularx}{\textwidth}{|l|X|}
        \hline
        \highlightbf{ABCFlyweight} & \\
        \hline
        \highlightbf{SharedFlyweight} & \\
        \hline
        \highlightbf{UnsharedFlyweight} & \\
        \hline
        \highlightbf{FlyweightFactory} & a class to generate and share 
                                         a pool of flyweight objects \\
        \hline
        \highlightbf{Client} & \\
        \hline
    \end{tabularx}    
    \caption{Flyweight Class Diagram.}
    \label{tab:flyweight-class-diagram} 
    \end{center}
\end{table}
\newpage
\toclessnparagraph{Pros and Cons}
\begin{table}[h!t]
    \begin{tabularx}{\textwidth}{|l|X|}
        \hline
        \textbf{Pros} & saving RAM while managing several similar 
                        objects \\
        \hline
        \textbf{Cons} & $\bullet$ CPU cycles might increase while 
                        RAM consumption decreases \\
                      & $\bullet$ code complexity and more overhead 
                        to understand the separation of states \\
        \hline      
    \end{tabularx}
\end{table}

%~ \toclessnparagraph{Relation with other Design Patterns}
 %~ \begin{table}[h!t]
    %~ \begin{tabularx}{\textwidth}{|l|X|}
        %~ \hline
        %~ \highlightbf{Composite} & shared leaf nodes of a \textit{Composite} 
                             %~ tree can be implemented as \textit{Flyweight}s \\
        %~ \hline
        %~ \highlightbf{Facade} & make a single object that represent an 
                          %~ entire subsystem  \\
        %~ \hline      
        %~ \highlightbf{Singleton} & it is resembled by a \textit{Flyweight} 
                             %~ in case the reduction goes so well that 
                             %~ the shared states of the objects can be 
                             %~ wrapped in a single flyweight. \\
                           %~ & The difference is that there should be 
                             %~ one Singleton instance, whereas there 
                             %~ can be multiple instances of flyweights; 
                             %~ moreover the Singleton can mutate, 
                             %~ flyweights can not.\\
        %~ \hline      
    %~ \end{tabularx}
%~ \end{table}
    
\newpage
%%% %%% PROXY %%% %%% %%%
\subsection{Proxy}
\begin{table}[h!t]
    \begin{tabularx}{\textwidth}{|l|X|}
        \hline
        \multirow{1}{*}{\textbf{What for}} & \highlightit{substitute 
                                             or placeholder for another 
                                             object}, \highlightit{controlling 
                                             access} to it giving the 
                                             opportunity to make actions
                                             both before and after 
                                             using the object\\
        \hline
        \multirow{1}{*}{\textbf{How to}} & a proxy class with the same 
                                           interace of the original object 
                                           is responsible of accepting 
                                           request and creating real 
                                           service objects\\
        \hline
        \textbf{When needed} & $\bullet$ lazy initialization\\
                             & $\bullet$ access control\\
                             & $\bullet$ local execution of a remote 
                               service \\
                             & $\bullet$ loggin the requests history 
                               through a proxy\\
                             & $\bullet$ caching requests result, 
                               expecially useful when results are 
                               large\\
                             & $\bullet$ using references to keep 
                               track of the clients using a service, 
                               in order \\
                             & \quad to dismiss it when not needed\\
        \hline
    \end{tabularx}
\end{table}

\toclessnparagraph{Structure}
\vspace{-1.5cm}
\begin{table}[h!t]
    \begin{center}
    \begingroup
    \begin{figure}[H]
        \begin{center}
    
    \begin{tikzpicture}[show background grid]
        
        \begin{abstractclass}[text width=4cm]{ABCServiceInterface}{0,0}
            \operationtt{+ operation()}
        \end{abstractclass}
        \begin{class}[text width=4cm]{Proxy}{-2.5,-4}
            \implement{ABCServiceInterface}
            \attributett{- realService: Service}
            \operationtt{+ Proxy(s: Service)}
            \operationtt{+ checkAccess()}
            \operationtt{+ operation()}
        \end{class}
        \begin{class}[text width=3cm]{Service}{2.5,-4}
            \implement{ABCServiceInterface}
            \attributett{...}
            \operationtt{+ operation()}
        \end{class}
        
        \begin{class}[text width=3cm]{Client}{-5,-0.2}
        \end{class}
        
    \end{tikzpicture}
    
\end{center}
       
    \end{figure}
    \endgroup
    \begin{tabularx}{\textwidth}{|l|X|}
        \hline
        \highlightbf{ABCServiceInterface} & the interface of the service 
                                            which requires a proxy, 
                                            that must follow this 
                                            interface to disguise 
                                            itself as the service.\\
        \hline
        \highlightbf{Service} & class providing some useful business logic.\\
        \hline
        \highlightbf{Proxy} & class containing a reference to a service 
                              object, to which it passes the request 
                              after some preliminary processing and 
                              of which it manages the lifecycle.\\
        \hline
        \highlightbf{Client} & works with both the service and the 
                               proxy, if they follow the same 
                               interface.\\
        \hline
    \end{tabularx}    
    \caption{Proxy Class Diagram.}
    \label{tab:proxy-class-diagram} 
    \end{center}    
\end{table}

\newpage
\toclessnparagraph{Pros and Cons}
\begin{table}[h!t]
    \begin{tabularx}{\textwidth}{|l|X|}
        \hline
        \textbf{Pros} & $\bullet$ control over service objects without 
                        any knoledge about their internals \\
                      & $\bullet$ lifecycle of the service object 
                        does not require management by the client \\
                      & $\bullet$ proxy can work while the service is 
                        initializing making programs faster \\
                      & $\bullet$ \highlightit{Open/Close Principle} 
                        respected since new proxies can be added without \\
                      & \quad changing the service or its clients\\
        \hline
        \textbf{Cons} & $\bullet$ new classes make code more complex \\
                      & $\bullet$ responses from service can be delayed \\
        \hline      
    \end{tabularx}
\end{table}
    
%~ \toclessnparagraph{Relation with other Design Patterns}
 %~ \begin{table}[h!t]
    %~ \begin{tabularx}{\textwidth}{|l|X|}
        %~ \hline
        %~ \highlightbf{Adapter} & $\bullet$ similarly, it garantees access 
                           %~ to an exising object; \\
                         %~ & $\bullet$ differently, the interface is 
                           %~ different than the adaptee one. \\
        %~ \hline
        %~ \highlightbf{Decorator} & $\bullet$ similarly to \textit{Proxy} it 
                             %~ guarantees access to a different object;\\
                           %~ & $\bullet$ differently, the object is accessed 
                             %~ through an enhanced interface. \\
        %~ \hline
        %~ \highlightbf{Facade} & it is similar to \textit{Proxy} but the 
                          %~ latter has the same interface of the object 
                          %~ and it is interchangeable; \\
        %~ \hline      
        %~ \highlightbf{Decorator} & $\bullet$ it also relies on \highlightit{composition principle}, and 
                             %~ one object delegates some work to another;  \\
                           %~ & $\bullet$ differently, \textit{Proxy} usually 
                             %~ manages the lifecycle of the object. \\
        %~ \hline      
    %~ \end{tabularx}
%~ \end{table}   
    
    
